%%%%%%%%%%%%%%%%%%%%%%%%%%%%%%%%%
%%\newpage
%\part{Patterson--Sullivan theory}
%%\section{Patterson--Sullivan theory}
%\label{partahlforsthurston}
%This part will be divided as follows: In Section \ref{sectionconformal} we recall the definition of quasiconformal measures, and we prove basic existence and non-existence results. In Section \ref{sectionahlforsthurston}, we prove Theorem \ref{theoremahlforsthurstongeneral} (Patterson--Sullivan theorem for groups of divergence type). In Section \ref{sectionGFmeasures}, we investigate the geometry of quasiconformal measures of geometrically finite groups, and we prove a generalization of the Global Measure Formula (Theorem \ref{theoremglobalmeasure}) as well as giving various necessary and/or sufficient conditions for the Patterson--Sullivan measure of a geometrically finite group to be doubling (\6\ref{subsectiondoubling}) or exact dimensional (\6\ref{subsectionexactdimensional}).
%
%Throughout the entire part, we fix $(X,\dist,\zero,b)$ as in \6\ref{standingassumptions2}, and a \emph{group} $G\leq\Isom(X)$.


%%%%%%%%%%%%%%%%%%%%%%%%%%%%%%%%%%
%\draftnewpage
\chapter{Quasiconformal measures of geometrically finite groups}
\label{sectionGFmeasures}
%%%%%%%%%%%%%%%%%%%%%%%%%%%%%%%%%%

In this chapter we investigate the $\delta$-quasiconformal measure or measures associated to a geometrically finite group. Note that since geometrically finite groups are of compact type (Theorem \ref{theoremGFcompact}), Theorem \ref{theorempattersonsullivangeneral} guarantees the existence of a $\delta$-quasiconformal measure $\mu$ on $\Lambda$. However, this measure is not necessarily unique (Corollary \ref{corollaryuniquenesscharacterization}); a sufficient condition for uniqueness is that $G$ is of divergence type (Theorem \ref{theoremahlforsthurstongeneral}). In Section \ref{subsectiondivergencetype}, we generalize a theorem of Dal'bo, Otal, and Peigne \cite[Th\'eor\`eme A]{DOP} which shows that ``most'' geometrically finite groups are of divergence type. In Sections \ref{subsectionglobalmeasure}-\ref{subsectionexactdimensional} we investigate the geometry of $\delta$-conformal measures; specifically, in Sections \ref{subsectionglobalmeasure}-\ref{subsectionglobalmeasureproof} we prove a generalization of the Global Measure Formula (Theorem \ref{theoremglobalmeasure}), in Sections \ref{subsectiondoubling} and \ref{subsectionexactdimensional} we investigate the questions of when the $\delta$-conformal measure of a geometrically finite group is doubling and exact dimensional, respectively.

\begin{standingassumptions}
\label{standingassumptionsGFmeasures}
In this chapter, we assume that
\begin{itemize}
\item[(I)] $X$ is regularly geodesic and strongly hyperbolic,
\item[(II)] $G\leq\Isom(X)$ is nonelementary and geometrically finite, and $\delta < \infty$.\Footnote{Note that by Corollary \ref{corollaryGF}(ii), we have $\delta < \infty$ if and only if $\delta_\bp < \infty$ for all $\bp\in P$.}
%\item[(III)] $\mu$ is a $\delta$-conformal measure on $\Lambda$.
\end{itemize}
Moreover, we fix a complete set of inequivalent parabolic points $P\subset\Lbp$, and for each $\bp\in P$ we write $\delta_\bp = \delta(G_\bp)$, and let $S_\bp\subset\EE_\bp$ be a $\bp$-bounded set satisfying (A)-(C) of Lemma \ref{lemmaboundedparabolic}. Finally, we choose a number $t_0 > 0$ large enough so that if
\begin{align*}
H_\bp &= H_{\bp,t_0} = \{ x\in X : \busemann_\bp(\zero,x) > t_0 \}\\
\scrH &= \{ g(H_\bp) : \bp\in P, g\in G\},
\end{align*}
then the collection $\scrH$ is disjoint (cf. Proof of Theorem \ref{theoremGFcompact}\text{(B3) \implies (A)}).
\end{standingassumptions}


%large number $t_0 > 0$ to be determined later (cf. Proposition \ref{propositionmetadecreasing}), and we let
%\begin{align*}
%H_\bp &= H_{\bp,t_0} = \{ x\in X : \busemann_\bp(\zero,x) > t_0 \}\\
%\scrH &= \{ g(H_\bp) : \bp\in P, g\in G\}.
%\end{align*}
%Recall (Proof of Theorem \ref{theoremGFcompact}(B3) \implies (A)) that for $t_0$ large enough, the collection $\scrH$ defined in this way satisfies the requirements of Definition \ref{definitionGF}.


\section{Sufficient conditions for divergence type}
\label{subsectiondivergencetype}

In the Standard Case, all geometrically finite groups are of divergence type \cite[Proposition 2]{Sullivan_entropy}; however, once one moves to the more general setting of pinched Hadamard manifolds, one has examples of geometrically finite groups of convergence type \cite[Th\'eor\`eme C]{DOP}. On the other hand, Proposition \ref{propositiondivtypedichotomy} shows that for every $\delta$-conformal measure $\mu$, $G$ is of divergence type if and only if $\mu(\Lambda\butnot\Lr) = 0$. Now by Theorem \ref{theoremGFcompact}, $\Lambda\butnot\Lr = \Lbp = G(P)$, so the condition $\mu(\Lambda\butnot\Lr) = 0$ is equivalent to the condition $\mu(P) = 0$. To summarize:

\begin{observation}
\label{observationnonatomic}
The following are equivalent:
\begin{itemize}
\item[(A)] $G$ is of divergence type.
\item[(B)] There exists a $\delta$-conformal measure $\mu$ on $\Lambda$ satisfying $\mu(P) = 0$.
\item[(C)] Every $\delta$-conformal measure $\mu$ on $\Lambda$ satisfies $\mu(P) = 0$.
\item[(D)] There exists a unique $\delta$-conformal measure $\mu$ on $\Lambda$, and it satisfies $\mu(P) = 0$.
\end{itemize}
In particular, every convex-cobounded group is of divergence type.
\end{observation}

It is of interest to ask for sufficient conditions which are not phrased in terms of measures. We have the following:

\begin{theorem}[Cf. {\cite[Proposition 2]{Sullivan_entropy}}, {\cite[Th\'eor\`eme A]{DOP}}]
\label{theoremGFdivergencetype}
If $\delta > \delta_\bp$ for all $\bp\in P$, then $G$ is of divergence type.
\end{theorem}
\begin{proof}
We will demonstrate (B) of Observation \ref{observationnonatomic}. Let $\mu$ be the measure constructed in the proof of Theorem \ref{theorempattersonsullivangeneral}, fix $\bp\in P$, and we will show that $\mu(\bp) = 0$. In what follows, we use the same notation as in the proof of Theorem \ref{theorempattersonsullivangeneral}. Since $G$ is strongly discrete, we can let $\rho$ be small enough so that $S_\rho = G(\zero)$. For any neighborhood $U$ of $\bp$, we have
\begin{equation}
\label{pattersonsullivanbound}
\mu(\bp) \leq \liminf_{s\searrow\delta} \mu_s(U)
= \liminf_{s\searrow\delta} \frac{1}{\Sigma_{s,k}} \sum_{x\in G(\zero)\cap U} k(x) e^{-s\dox x}.
\end{equation}

\begin{lemma}
\label{lemmaSbpasymp}
\[
\lb h(\zero) | x\rb_\zero \asymp_\plus 0 \all x\in S_\bp.
\]
\end{lemma}
\begin{subproof}
Since $S_\bp$ is $\bp$-bounded, Gromov's inequality implies that
\[
\lb h(\zero) | x\rb_\zero \wedge \lb h(\zero)|\bp\rb_\zero \asymp_\plus 0
\]
for all $h\in G_\bp$ and $x\in S_\bp$. Denote the implied constant by $\sigma$. For all $h\in G_\bp$ such that $\lb h(\zero)|\bp\rb_\zero > \sigma$, we have $\lb h(\zero) | x\rb_\zero \leq \sigma \all x\in S_\bp$. Since this applies to all but finitely many $h\in G_\bp$, (c) of Proposition \ref{propositionbasicidentities} completes the proof.
\end{subproof}

Let $T$ be a transversal of $G_\bp\backslash G$ such that $T(\zero) \subset S_\bp$. Then by Lemma \ref{lemmaSbpasymp},
\[
\dox{h(x)} \asymp_\plus \dogo h + \dox x \all h\in G_\bp \all x\in T(\zero).
\]
Thus for all $s > \delta$ and $V\subset X$,
\begin{equation}
\label{SigmaVasymp}
\begin{split}
\sum_{x\in G(\zero)\cap U} k(x) e^{-s\dox x}
&=_\pt \sum_{h\in G_\bp} \sum_{x\in hT(\zero)\cap U} k(e^{\dox x}) e^{-s\dox x}\\
&\asymp_\times \sum_{h\in G_\bp} \sum_{x\in T(\zero)\cap h^{-1}(U)} k(e^{\dogo h + \dox x}) e^{-s[\dogo h + \dox x]}.
\end{split}
\end{equation}
Now fix $0 < \epsilon < \delta - \delta_\bp$, and note that by \eqref{patterson},
\[
k(R) \leq k(\lambda R) \lesssim_{\times,\epsilon} \lambda^\epsilon k(R) \all \lambda > 1 \all R \geq 1.
\]
Thus setting $V = U$ in \eqref{SigmaVasymp} gives
\[
\sum_{x\in G(\zero)\cap U} k(x) e^{-s\dox x} \lesssim_{\times,\epsilon}  \sum_{\substack{h\in G_\bp \\ h(S_\bp)\cap U\neq\smallemptyset}} e^{-(s - \epsilon)\dogo h} \sum_{x\in T(\zero)} k(x)e^{-s\dox x},
\]
while setting $V = X$ gives
\begin{align*}
\Sigma_{s,k} = \sum_{x\in G(\zero)} k(x) e^{-s\dox x} \gtrsim_\times \sum_{h\in G_\bp} e^{-s\dogo h} \sum_{x\in T(\zero)} k(x)e^{-s\dox x}.
\end{align*}
Dividing these inequalities and combining with \eqref{pattersonsullivanbound} gives
\[
\mu(\bp)
\lesssim_{\times,\epsilon} \liminf_{s\searrow\delta} \frac{1}{\Sigma_s(G_\bp)} \sum_{\substack{h\in G_\bp \\ h(S_\bp)\cap U \neq\smallemptyset}} e^{-(s - \epsilon)\dogo h}
= \frac{1}{\Sigma_\delta(G_\bp)} \sum_{\substack{h\in G_\bp \\ h(S_\bp)\cap U \neq\smallemptyset}} e^{-(\delta - \epsilon)\dogo h} ~.
\]
Note that the right hand series converges since $\delta - \epsilon > \delta_\bp$ by construction. As the neighborhood $U$ shrinks, the series converges to zero. This completes the proof.
\end{proof}

Combining Theorem \ref{theoremGFdivergencetype} with Proposition \ref{propositiondeltagtrdeltap} gives the following immediate corollary:

\begin{corollary}
If for all $\bp\in P$, $G_\bp$ is of divergence type, then $G$ is of divergence type.
\end{corollary}
Thus in some sense divergence type can be ``checked locally'' just like the properties of finite generation and finite Poincar\'e exponent (cf. Corollary \ref{corollaryGF}).
\begin{corollary}
Every convex-cobounded group is of divergence type.
\end{corollary}

\begin{remark}
It is somewhat awkward that it seems to be difficult or impossible to prove Theorem \ref{theoremGFdivergencetype} via any of the equivalent conditions of Observation \ref{observationnonatomic} other than (B). Specifically, the fact that the above argument works for the measure constructed in Theorem \ref{theorempattersonsullivangeneral} (the ``Patterson--Sullivan measure'') but not for other $\delta$-conformal measures seems rather asymmetric. However, after some thought one realizes that it would be impossible for a proof along similar lines to work for every $\delta$-conformal measure. This is because the above proof shows that the Patterson--Sullivan measure $\mu$ satisfies
\begin{equation}
\label{pattersonsullivanproperty}
\text{$\mu(\bp) = 0$ for all $\bp\in P$ satisfying $\delta > \delta_\bp$},
\end{equation}
but there are geometrically finite groups for which \eqref{pattersonsullivanproperty} does not hold for all $\delta$-conformal measures $\mu$. Specifically, one may construct geometrically finite groups of convergence type (cf. \cite[Th\'eor\`eme C]{DOP}) such that $\delta_\bp < \delta$ for some $\bp\in P$; the following proposition shows that there exists a $\delta$-conformal measure for which \eqref{pattersonsullivanproperty} fails:
\end{remark}

\begin{proposition}
\label{propositionGFconvergencetype}
If $G$ is of convergence type, then for each $\bp\in P$ there exists a $\delta$-conformal measure supported on $G(\bp)$.
\end{proposition}
\begin{proof}
Let
\[
\mu = \sum_{g(\bp)\in G(\bp)} [g'(\bp)]^\delta \delta_{g(\bp)};
\]
clearly $\mu$ is a $\delta$-conformal measure, but we may have $\mu(\del X) = \infty$. To prove that this is not the case, as before we let $T$ be a transversal of $G_\bp\backslash G$ such that $T(\zero)\subset S_\bp$. Then
\[
\mu(\del X) = \sum_{g(\bp)\in G(\bp)} [g'(\bp)]^\delta = \sum_{g\in T^{-1}} [g'(\bp)]^\delta \asymp_\times \sum_{g\in T^{-1}} e^{-\delta\dogo g} \leq \Sigma_\delta(G) < \infty.
\]
\end{proof}
Proposition \ref{propositionGFconvergencetype} yields the following characterization of when there exists a unique $\delta$-conformal measure:
\begin{corollary}
\label{corollaryuniquenesscharacterization}
The following are equivalent:
\begin{itemize}
\item[(A)] There exists a unique $\delta$-conformal measure on $\Lambda$.
\item[(B)] Either $G$ is of divergence type, or $\#(P) = 1$.
\end{itemize}
\end{corollary}

\section{The global measure formula}
\label{subsectionglobalmeasure}
In this section and the next, we fix a $\delta$-quasiconformal measure $\mu$, and ask the following geometrical question: Given $\eta\in\Lambda$ and $r > 0$, can we estimate $\mu(B(\eta,r))$? If $G$ is convex-cobounded, then we can show that $\mu$ is Ahlfors $\delta$-regular (Corollary \ref{corollaryCCBregular}), but in general the measure $\mu(B(\eta,r))$ will depend on the point $\eta$, in a manner described by the \emph{global measure formula}. To describe the global measure formula, we need to introduce some notation:

\begin{notation}
\label{notationmetat}
Given $\xi = g(\bp)\in\Lbp$, let $t_\xi > 0$ be the unique number such that
\[
H_\xi = H_{\xi,t_\xi} = g(H_\bp) = g(H_{\bp,t_0}),
\]
i.e. $t_\xi = t_0 + \busemann_\xi(\zero,g(\zero))$. (Note that $t_\bp = t_0$ for all $\bp\in P$.) Fix $\theta > 0$ large to be determined below (cf. Proposition \ref{propositionmetadecreasing}). For each $\eta\in\Lambda$ and $t > 0$, let $\eta_t = \geo\zero\eta_t$, and write
\begin{equation}
\label{metat}
m(\eta,t) = \begin{cases}
e^{-\delta t} & \eta_t\notin\bigcup(\scrH) \\
e^{-\delta t_\xi} [\II_\bp(e^{t - t_\xi - \theta}) + \mu(\bp)] & \eta_t\in H_\xi \text{ and }t \leq \lb \xi|\eta\rb_\zero\\
e^{-\delta (2\lb \xi|\eta\rb_\zero - t_\xi)} \NN_\bp(e^{2\lb \xi|\eta\rb_\zero - t - t_\xi - \theta}) & \eta_t\in H_\xi \text{ and } t > \lb \xi|\eta\rb_\zero
\end{cases}
\end{equation}
(cf. Figure \ref{figuremetat}.) Here we use the notation
\begin{align*}
\II_\bp(R) &= \sum_{\substack{h\in G_\bp \\ \|h\|_\bp \geq R}} \|h\|_\bp^{-2\delta}\\
\NN_\bp(R) &= \NN_{\EE_\bp,G_\bp}(R) = \#\{h\in G_\bp : \|h\|_\bp \leq R\}
\end{align*}
where
\[
\|h\|_\bp = \Dist_\bp(\zero,h(\zero)) = e^{(1/2)\dogo h}
\all h\in G_\bp.
\]
\end{notation}



\begin{figure}
%\begin{tikzpicture}[line cap=round,line join=round,>=triangle 45,x=1.0cm,y=1.0cm]
\begin{tikzpicture}[line cap=round,line join=round,>=triangle 45,scale=0.45]
\clip(-0.8063966880341881,-10.0638437660102) rectangle (14.110082748891738,6.841499595839061);
\draw (0.0,4.0)-- (1.0,4.0);
\draw (2.0,5.0)-- (3.0,4.0);
\draw (2.0,5.0)-- (1.0,4.0);
\draw (3.0,4.0)-- (5.0,4.0);
\draw (7.0,6.0)-- (5.0,4.0);
\draw (7.0,6.0)-- (9.0,4.0);
\draw (9.0,4.0)-- (12.0,4.0);
\draw (0.0,3.0)-- (1.0,2.0);
\draw [dash pattern=on 4pt off 4pt] (1.0,2.0)-- (2.0,1.0);
\draw [dash pattern=on 4pt off 4pt] (2.0,1.0)-- (3.0,0.0);
\draw (3.0,0.0)-- (5.0,-2.0);
\draw [dash pattern=on 4pt off 4pt] (5.0,-2.0)-- (9.0,-6.0);
\draw (9.0,-6.0)-- (12.0,-9.0);
\draw (1.0,2.0)-- (2.009795868802983,1.4948904390904918);
\draw (2.009795868802983,1.4948904390904918)-- (3.0,0.0);
\draw (5.0,-2.0)-- (7.006776519282634,-4.79037428534092);
\draw (7.006776519282634,-4.79037428534092)-- (9.0,-6.0);
\draw (10.285344431731245,5.388099035318084) node[anchor=north west] {$b(\eta,t)$};
\draw (10.0,-6.0) node[anchor=north west] {$\log m(\eta,t)$};
\end{tikzpicture}
\caption[Cusp excursion and ball measure functions]{A possible (approximate) graph of the functions $t\mapsto b(\eta,t)$ and $t\mapsto \log m(\eta,t)$ (cf. \eqref{metat} and \eqref{betat}). The graph indicates that there are at least two inequivalent parabolic points $\bp_1,\bp_2\in P$, which satisfy $\NN_{\bp_i}(R) \asymp_\times R^{2\delta} \II_{\bp_i}(R) \asymp_\times R^{k_i}$ for some $k_1 < 2\delta < k_2$. The dotted line in the second graph is just the line $y = -\delta t$.\newline
Note the relation between the two graphs, which may be either direct or inverted depending on the functions $\NN_\bp$. Specifically, the relation is direct for the first cusp but inverted for the second cusp.}
%The function $t\mapsto b(\eta,t)$ is piecewise linear, and its derivative takes values in $\{0,1,-1\}$. Intuitively, it tells you the ``distance penetrated into the cusp'' (cf. Observation \ref{observationbetatinterpretation}); when it is increasing, you are moving closer to the cusp, and when it is decreasing, you are moving farther away.\newline
%The function $t\mapsto m(\eta,t)$ can behave somewhat erratically, but it is always (approximately) decreasing (Proposition \ref{propositionmetadecreasing}).
\label{figuremetat}
\end{figure}


\begin{theorem}[Global measure formula; cf. {\cite[Theorem 2]{StratmannVelani}} and {\cite[Th\'eor\`eme 3.2]{Schapira}}]
\label{theoremglobalmeasure}
For all $\eta\in\Lambda$ and $t > 0$,
\begin{equation}
\label{globalmeasureformula}
m(\eta,t + \sigma) \lesssim_\times \mu(B(\eta,e^{-t})) \lesssim_\times m(\eta,t - \sigma),
\end{equation}
where $\sigma > 0$ is independent of $\eta$ and $t$ (but may depend on $\theta$).
\end{theorem}

\begin{corollary}
\label{corollaryCCBregular}
If $G$ is convex-cobounded, then
\begin{equation}
\label{ahlforsregular}
\mu(B(\eta,r)) \asymp_\times r^\delta \all \eta\in\Lambda \all 0 < r \leq 1,
\end{equation}
i.e. $\mu$ is Ahlfors $\delta$-regular.
\end{corollary}
\begin{proof}
If $G$ is convex-cobounded then $\scrH = \emptyset$, so $m(\eta,t) = e^{-\delta t}\all\eta,t$, and thus \eqref{globalmeasureformula} reduces to \eqref{ahlforsregular}.
\end{proof}

\begin{remark}
Corollary \ref{corollaryCCBregular} can be deduced directly from Lemma \ref{lemmaCCBglobalmeasure} below.
\end{remark}

We will prove Theorem \ref{theoremglobalmeasure} in the next section. For now, we investigate more closely the function $t\mapsto m(\eta,t)$ defined by \eqref{metat}. The main result of this section is the following proposition, which will be used in the proof of Theorem \ref{theoremglobalmeasure}:

\begin{proposition}
\label{propositionmetadecreasing}
If $\theta$ is chosen sufficiently large, then for all $\eta\in\Lambda$ and $0 < t_1 < t_2$,
\begin{equation}
\label{metadecreasing}
m(\eta,t_2) \lesssim_{\times,\theta} m(\eta,t_1).
\end{equation}
\end{proposition}

The proof of Proposition \ref{propositionmetadecreasing} itself requires several lemmas.

\begin{lemma}
\label{lemmabetatinterpretation}
Fix $\xi,\eta\in \del X$ and $t > 0$, and let $x = \eta_t$. Then
\begin{equation}
\label{betatinterpretation}
\busemann_\xi(\zero,x) \asymp_\plus t \wedge (2\lb\xi|\eta\rb_\zero - t).
\end{equation}
\end{lemma}
\begin{proof}
Since $\lb \zero|\eta\rb_x = 0$, Gromov's inequality gives $\lb \zero|\xi\rb_x \wedge \lb \xi|\eta\rb_x \asymp_\plus 0$.

\begin{itemize}
\item[Case 1:] $\lb \zero|\xi\rb_x \asymp_\plus \zero$. In this case, by (h) of Proposition \ref{propositionbasicidentities},
\[
\busemann_\xi(\zero,x) = - \busemann_\xi(x,\zero) = -[2\lb \zero|\xi\rb_x - \dox x] \asymp_\plus \dox x = t,
\]
while (g) of Proposition \ref{propositionbasicidentities} gives
\[
\lb \xi|\eta\rb_\zero = \lb \xi|\eta\rb_x + \frac12[\busemann_\xi(\zero,x) + \busemann_\eta(\zero,x)] \gtrsim_\plus \frac12[t + t] = t;
\]
thus $\busemann_\xi(\zero,x) \asymp_\plus t \asymp_\plus t\wedge(2\lb \xi|\eta\rb_\zero - t)$.
\item[Case 2:] $\lb \xi|\eta\rb_x \asymp_\plus \zero$. In this case, (g) of Proposition \ref{propositionbasicidentities} gives
\[
\lb \xi|\eta\rb_\zero \asymp_\plus \frac12[\busemann_\xi(\zero,x) + \busemann_\eta(\zero,x)] = \frac12[\busemann_\xi(\zero,x) + t] \lesssim_\plus \frac12[t + t] = t;
\]
thus $\busemann_\xi(\zero,x) \asymp_\plus 2\lb \xi|\eta\rb_\zero - t \asymp_\plus t\wedge(2\lb \xi|\eta\rb_\zero - t)$.
\end{itemize}
\end{proof}

\begin{corollary}
\label{corollarybetat}
The function
\begin{equation}
\label{betat}
b(\eta,t) = \begin{cases}
0 & \eta_t\notin\bigcup(\scrH) \\
t\wedge(2\lb \xi|\eta\rb_\zero - t) - t_\xi & \eta_t\in H_\xi
\end{cases}
\end{equation}
satisfies
\begin{equation}
\label{betatdoubling}
b(\eta,t + \tau) \asymp_{\plus,\tau} b(\eta,t - \tau).
\end{equation}
\end{corollary}
\begin{proof}
Indeed, by Lemma \ref{lemmabetatinterpretation},
\begin{align*}
b(\eta,t) &\asymp_\plus \begin{cases}
0 & \eta_t\notin\bigcup(\scrH) \\
\busemann_\xi(\zero,\eta_t) - t_\xi & \eta_t\in H_\xi
\end{cases}\\
&=_\pt 0\vee \max_{\xi\in\Lbp} (\busemann_\xi(\zero,\eta_t) - t_\xi).
\end{align*}
The right hand side is $1$-Lipschitz continuous with respect to $t$, which demonstrates \eqref{betatdoubling}.
\end{proof}



\begin{lemma}
\label{lemmatop}
For all $\xi\in G(\bp) \subset \Lbp$, $\bp\in P$, there exists $g\in G$ such that
\begin{equation}
\label{top}
\text{$\xi = g(\bp)$, $\dogo g \asymp_\plus t_\xi$, and $\{\eta\in\del X  : \geo\zero\eta\cap H_\xi\neq\emptyset \} \subset \Shad(g(\zero),\sigma)$,}
\end{equation}
where $\sigma > 0$ is independent of $\xi$.
\end{lemma}
\begin{proof}
Write $\xi = g(\bp)$ for some $g\in G$. Since $x := \xi_{t_\xi} \in \del H_\xi$, Lemma \ref{lemmaboundedparabolic}(D) shows that
\[
\dist(g^{-1}(x),h(\zero)) \asymp_\plus 0
\]
for some $h\in G_\bp$. We claim that $gh$ is the desired isometry. Clearly $\dogo{gh} \asymp_\plus \dox x = t_\xi$. Fix $\eta\in\del X$ such that $\geo\zero\eta\cap H_\xi\neq\emptyset$, say $\eta_t\in H_\xi$. By Lemma \ref{lemmabetatinterpretation}, we have
\[
\dox x = t_\xi < \busemann_\xi(\zero,\eta_t) \asymp_\plus t\wedge(2\lb \xi|\eta\rb_\zero - t) \leq \lb \xi|\eta\rb_\zero \leq \lb x|\eta\rb_\zero,
\]
i.e. $\eta\in \Shad(x,\sigma) \subset \Shad(g(\zero),\sigma + \tau)$ for some $\sigma,\tau > 0$.
\end{proof}

\begin{proof}[Proof of Proposition \ref{propositionmetadecreasing}]
Fix $\eta\in\Lambda$ and $0 < t_1 < t_2$.
\begin{itemize}
\item[Case 1:] $\eta_{t_1},\eta_{t_2}\in H_\xi$ for some $\xi = g(\bp) \in \Lbp$, $g$ satisfying \eqref{top}. In this case, \eqref{metadecreasing} follows immediately from \eqref{metat} unless $t_1 \leq \lb \xi|\eta\rb_\zero < t_2$. If the latter holds, then
\begin{align*}
m(\eta,t_1) &\geq \lim_{t\nearrow\lb \xi|\eta\rb_\zero}m(\eta,t) = e^{-\delta t_\xi} [\II_\bp(e^{\lb \xi|\eta\rb_\zero - t_\xi - \theta}) + \mu(\bp)]\\
m(\eta,t_2) &\leq \lim_{t\searrow\lb \xi|\eta\rb_\zero}m(\eta,t) = e^{-\delta(2\lb \xi|\eta\rb_\zero - t_\xi)} \NN_\bp(e^{\lb \xi|\eta\rb_\zero - t_\xi - \theta}).
\end{align*}
Consequently, to demonstrate \eqref{metadecreasing} it suffices to show that
\begin{equation}
\label{ETSmetadecreasing}
\NN_\bp(e^t) \lesssim_{\times,\theta} e^{2\delta t}\II_\bp(e^t),
\end{equation}
where $t := \lb \xi|\eta\rb_\zero - t_\xi - \theta > 0$.

To demonstrate \eqref{ETSmetadecreasing}, let $\zeta = g^{-1}(\eta)\in\Lambda$. We have
\[
\lb \bp|\zeta\rb_\zero = \lb \xi|\eta\rb_{g(\zero)} \asymp_\plus \lb \xi|\eta\rb_\zero - \dogo g \asymp_\plus \lb \xi|\eta\rb_\zero - t_\xi = t + \theta
\]
and thus
\[
\Dist_\bp(\zero,\zeta) \asymp_\times e^{t + \theta}.
\]
Since $\bp$ is a bounded parabolic point, there exists $h_\zeta\in G_\bp$ such that $\Dist_\bp(h_\zeta(\zero),\zeta) \lesssim_\times 1$. Denoting all implied constants by $C$, we have
\begin{align*}
C^{-1} e^{t + \theta} - C \leq \Dist_\bp(\zero,\zeta) - \Dist_\bp(h_\zeta(\zero),\zeta) 
&\leq \|h_\zeta\|_\bp\\ 
&\leq \Dist_\bp(\zero,\zeta) + \Dist_\bp(h_\zeta(\zero),\zeta) \leq C e^{t + \theta} + C.
\end{align*}
Choosing $\theta \geq \log(4C)$, we have
\[
2 e^t \leq \|h_\zeta\|_\bp \leq 2C e^{t + \theta} \text{ unless $e^{t + \theta} \leq 2C^2$}.
\]
If $2 e^t \leq \|h_\zeta\|_\bp \leq 2C e^{t + \theta}$, then for all $h\in G_\bp$ satisfying $\|h\|_\bp\leq e^t$ we have $e^t \leq \|h_\zeta h\|_\bp \lesssim_{\times,\theta} e^t$; it follows that
\[
\II_\bp(e^t) \geq \sum_{h\in G_\bp} \|h_\zeta h\|_\bp^{-2\delta} \asymp_{\times,\theta} e^{-2\delta t} \NN_\bp(e^t),
\]
thus demonstrating \eqref{ETSmetadecreasing}. On the other hand, if $e^{t + \theta} \leq 2C^2$, then both sides of \eqref{ETSmetadecreasing} are bounded from above and below independent of $t$.
\item[Case 2:] No such $\xi$ exists. In this case, for each $i$ write $\eta_i\in H_{\xi_i}$ for some $\xi_i = g_i(\bp_i)\in\Lbp$ if such a $\xi_i$ exists. If $\xi_1$ exists, let $s_1 > t_1$ be the smallest number such that $\eta_{s_1}\in\del H_{\xi_1}$, and if $\xi_2$ exists, let $s_2 < t_2$ be the largest number such that $\eta_{s_2}\in \del H_{\xi_2}$. If $\xi_i$ does not exist, let $s_i = t_i$. Then $t_1 \leq s_1 \leq s_2 \leq t_2$. Since $m(\eta,s_i) = e^{-\delta s_i}$, we have $m(\eta,s_2) \leq m(\eta,s_1)$, so to complete the proof it suffices to show that
\begin{align*}
m(\eta,s_1) &\lesssim_{\times,\theta} m(\eta,t_1) \text{ and}\\
m(\eta,s_2) &\gtrsim_{\times,\theta} m(\eta,t_2).
\end{align*}
By Case 1, it suffices to show that
\begin{align*}
m(\eta,s_1) &\lesssim_\times \lim_{t\nearrow s_1}m(\eta,t) \text{ if $\xi_1$ exists, and}\\
m(\eta,s_2) &\gtrsim_\times \lim_{t\searrow s_2} m(\eta,t) \text{ if $\xi_2$ exists.}
\end{align*}
Comparing with \eqref{metat}, we see that the desired formulas are
\begin{align*}
e^{-\delta s_1} &\lesssim_\times e^{-\delta (2\lb \xi|\eta\rb_\zero - t_{\xi_1})} \NN_\bp(e^{2\lb \xi_1|\eta\rb_\zero - s_1 - t_{\xi_1}})\\
e^{-\delta s_2} &\gtrsim_\times  e^{-\delta t_{\xi_2}}[\II_\bp(e^{s_2 - t_{\xi_2}}) + \mu(\bp)],
\end{align*}
which follow upon observing that the definitions of $s_1$ and $s_2$ imply that $s_1 \asymp_\plus 2\lb \xi|\eta\rb_\zero - t_{\xi_1}$ and $s_2 \asymp_\plus t_{\xi_2}$ (cf. Lemma \ref{lemmabetatinterpretation}).
\end{itemize}
\end{proof}

%\section{Proof of Theorem \ref{theoremglobalmeasure}}
\section{Proof of the global measure formula}
\label{subsectionglobalmeasureproof}
Although we have finished the proof of Proposition \ref{propositionmetadecreasing}, we still need a few lemmas before we can begin the proof of Theorem \ref{theoremglobalmeasure}. Throughout these lemmas, we fix $\bp\in P$, and let
\[
R_\bp = \sup_{x\in S_\bp} \Dist_\bp(\zero,x) < \infty.
\]
Here $S_\bp\subset\EE_\bp$ is a $\bp$-bounded set satisfying $\Lambda\butnot\{\bp\}\subset G_\bp(S_\bp)$, as in Standing Assumptions \ref{standingassumptionsGFmeasures}.

\begin{lemma}
\label{lemmaglobalmeasure1}
For all $A\subset G_\bp$,
\begin{equation}
\label{ASmeasure}
\mu\left(\bigcup_{h\in A}h(S_\bp)\right) \asymp_\times \sum_{h\in A} e^{-\delta\dogo h} = \sum_{h\in A} \|h\|_\bp^{-2\delta}.
\end{equation}
\end{lemma}
\begin{proof}
As the equality follows from Observation \ref{observationeuclideanparabolicasymp}, we proceed to demonstrate the asymptotic. By Lemma \ref{lemmaSbpasymp}, there exists $\sigma > 0$ such that $S_\bp \subset \Shad_{h^{-1}(\zero)}(\zero,\sigma)$ for all $h\in G_\bp$. Then by the Bounded Distortion Lemma \ref{lemmaboundeddistortion},
\[
\mu(h(S_\bp)) = \int_{S_\bp} (\overline h')^\delta \;\dee \mu \asymp_{\times,\sigma} e^{-\delta\dogo h} \mu(S_\bp) \asymp_\times e^{-\delta\dogo h}.
\]
(In the last asymptotic, we have used the fact that $\mu(S_\bp) > 0$, which follows from the fact that $\Lambda\butnot\{\bp\} \subset G_\bp(S_\bp)$ together with the fact that $\mu$ is not a pointmass (Corollary \ref{corollaryquasiconformalpointmass}).) Combining with the subadditivity of $\mu$ gives the $\lesssim$ direction of the first asymptotic of \eqref{ASmeasure}. To get the $\gtrsim$ direction, we observe that since $S_\bp$ is $\bp$-bounded, the strong discreteness of $G_\bp$ implies that $S_\bp\cap h(S_\bp)\neq\emptyset$ for only finitely many $h\in G_\bp$; it follows that the function $\eta\mapsto\#\{h\in G_\bp : \eta\in h(S_\bp)\}$ is bounded, and thus
\begin{align*}
\mu\left(\bigcup_{h\in A}h(S_\bp)\right) 
&\asymp_\times \int \#\{h\in G_\bp : \eta\in h(S_\bp)\} \;\dee\mu(\eta)\\
&= \sum_{h\in A}\mu(h(S_\bp))\\
&\asymp_\times \sum_{h\in A} e^{-\delta\dogo h} ~.
\end{align*}
\end{proof}


\begin{corollary}
\label{corollaryglobalmeasure11}
For all $r > 0$,
\begin{equation}
\label{muBbprbounds}
\II_\bp\left(\frac2r\right) \lesssim_\times \mu\big(B(\bp,r)\butnot\{\bp\}\big) \lesssim_\times \II_\bp\left(\frac1{2r}\right)
\end{equation}
\end{corollary}
\begin{proof}
Since
\[
\bigcup_{\substack{h\in G_\bp \\ \|h\|_\bp \geq R + R_\bp}} h(S_\bp)
\subset B(\bp,1/R)\butnot\{\bp\} = \EE_\bp \butnot B_\bp(\zero,R)
\subset \bigcup_{\substack{h\in G_\bp \\ \|h\|_\bp \geq R - R_\bp}} h(S_\bp),
\]
Lemma \ref{lemmaglobalmeasure1} gives
\[
\II_\bp\left(\frac1r + R_\bp\right) \lesssim_\times \mu(B(\bp,r)) \lesssim_\times \II_\bp\left(\frac1r - R_\bp\right),
\]
thus proving the lemma if $r \leq 1/(2R_\bp)$. But when $r > 1/(2 R_\bp)$, all terms of \eqref{muBbprbounds} are bounded from above and below independent of $r$.
\end{proof}

Adding $\mu(\bp)$ to all sides of \eqref{muBbprbounds} gives
\begin{equation}
\label{muBbprbounds2}
\II_\bp\left(\frac2r\right) + \mu(\bp) \lesssim_\times \mu(B(\bp,r)) \lesssim_\times \II_\bp\left(\frac1{2r}\right) + \mu(\bp).
\end{equation}

\begin{figure}
\begin{center}
\begin{tabular}{ll}
%\begin{tikzpicture}[line cap=round,line join=round,>=triangle 45,x=1.0cm,y=1.0cm]
\begin{tikzpicture}[line cap=round,line join=round,>=triangle 45,scale=.81]
%\clip(-4.3527228989498425,-4.317699333578228) rectangle (4.5670505753079516,4.22893620649748);
\clip(-2.7,-2.7) rectangle (4.3,4.3);
\draw(0.0,0.0) circle (4.0cm);
\draw(2.42888,0.1875) circle (1.551112790229163cm);
\draw (0.0,0.0)-- (4.0,0.0);
\draw(3.212205575454762,0.0) circle (0.3760152197599791cm);
%\draw (0.0,0.0)-- (3.9725001940685334,0.4682330703030985);
\draw [dash pattern=on 3pt off 3pt] (0.0,0.0)-- (3.9725001940685334,0.4682330703030985);
%\draw (0.0,0.0)-- (3.9725001940685334,-0.4682330703030985);
\draw [dash pattern=on 3pt off 3pt] (0.0,0.0)-- (3.9725001940685334,-0.4682330703030985);
\draw [shift={(0.0,0.0)},line width=1.2000000000000002pt]  plot[domain=-0.11732726409329874:0.11732726409329915,variable=\t]({1.0*4.0*cos(\t r)+-0.0*4.0*sin(\t r)},{0.0*4.0*cos(\t r)+1.0*4.0*sin(\t r)});
\begin{scriptsize}
\draw [fill=black] (0.0,0.0) circle (1pt);
\draw[color=black] (-0.2491340591615706,-0.0014223831473892945) node {$\zero$};
\draw [fill=black] (4.0,0.0) circle (1pt);
\draw[color=black] (4.211681114182143,0.0) node {$\eta$};
\draw [fill=black] (3.98,0.3) circle (1pt);
\draw[color=black] (4.17,0.4) node {$\xi$};
\draw [fill=black] (3.384151714048233,0.0) circle (1pt);
\draw[color=black] (3.2344150960861695,-0.1380274618110051) node {$\eta_t$};
\end{scriptsize}
\end{tikzpicture}
%\begin{tikzpicture}[line cap=round,line join=round,>=triangle 45,x=1.0cm,y=1.0cm]
\begin{tikzpicture}[line cap=round,line join=round,>=triangle 45,scale=.85]
\clip(-4.152066304477812,-0.497280356328431) rectangle (4.307443018392998,5.909099879057296);
\draw (-3.5,0.0)-- (3.5,0.0);
\draw (-3.5,1.5)-- (3.5,1.5);
\draw [shift={(0.40000000000000036,0.0)}] plot[domain=0.3996780691889681:3.141592653589793,variable=\t]({1.0*2.518491612056709*cos(\t r)+-0.0*2.518491612056709*sin(\t r)},{0.0*2.518491612056709*cos(\t r)+1.0*2.518491612056709*sin(\t r)});
\draw (2.4230669786242873,1.5)-- (2.1163456022508336,1.2434075389985795);
\draw (-2.1184916120567086,0.0)-- (-2.12,5.4);
\draw (2.1163456022508336,1.2434075389985795)-- (2.12,5.4);
\begin{scriptsize}
\draw [fill=black] (2.72,0.98) circle (1pt);
\draw[color=black] (2.9026704848505824,1.0) node {$\zero$};
\draw [fill=black] (-2.1184916120567086,0.0) circle (1pt);
\draw[color=black] (-2.284181946690644,-0.1267908969326299) node {$\eta$};
\draw [fill=black] (2.4230669786242873,1.5) circle (1pt);
\draw[color=black] (2.58,1.68) node {$\eta_{t_\xi}$};
\draw [fill=black] (2.1163456022508336,1.2434075389985795) circle (1pt);
\draw[color=black] (1.8073209579773947,1.092736906911882) node {$g(\zero)$};
%\draw [fill=black] (-0.13843049435662466,5.4459880548125446) circle (1pt);
\draw[color=black] (-0.1,5.662106906126763) node {$\xi=\infty$};
\draw [fill=black] (1.6640939374267694,2.1782714517159905) circle (1pt);
\draw[color=black] (1.420712647267375,2.1424570418666518) node {$\eta_t$};
\draw [dash pattern=on 1pt off 1pt] (1.6640939374267694,2.1782714517159905)-- (2.1173124637830534,2.3431371337242592);
\draw[color=black] (-3.084181946690644,1.867908969326299) node {$H_{\xi}$};
\end{scriptsize}
\end{tikzpicture}
\end{tabular}
\end{center}

\caption[Estimating measures of balls via information ``at infinity'']{Cusp excursion in the ball model (left) and upper half-space model (right). Since $\xi = g(\bp)\in B(\eta,e^{-t})$, our estimate of $\mu(B(\eta,e^{-t}))$ is based on the function $\II_\bp$, which captures information ``at infinity'' about the cusp $\bp$. In the right-hand picture, the measure of $B(\eta,e^{-t})$ can be estimated by considering the measure from the perspective of $g(\zero)$ of a small ball around $\xi$.}
\label{figurecuspright}
\end{figure}


\begin{corollary}[Cf. Figure \ref{figurecuspright}]
\label{corollaryglobalmeasure12}
Fix $\eta\in\Lambda$ and $t > 0$ such that $\eta_t\in H_\xi$ for some $\xi = g(\bp)\in\Lbp$ satisfying $t \leq \lb \xi|\eta\rb_\zero - \log(2)$. Then
\[
e^{-\delta t_\xi} [\II_\bp(e^{t - t_\xi + \sigma}) + \mu(\bp)] \lesssim_\times \mu\big(B(\eta,e^{-t})\big) \lesssim_\times e^{-\delta t_\xi} [\II_\bp(e^{t - t_\xi - \sigma}) + \mu(\bp)],
\]
where $\sigma > 0$ is independent of $\eta$ and $t$.
\end{corollary}
\begin{proof}
The inequality $\lb \xi|\eta\rb_\zero \geq t + \log(2)$ implies that
\[
B(\xi,e^{-t}/2) \subset B(\eta,e^{-t}) \subset B(\xi,2e^{-t}).
\]
Without loss of generality suppose that $g$ satisfies \eqref{top}. Since $t > t_\xi$, \eqref{ShadcontainsB} guarantees that $B(\xi,2e^{-t}) \subset \Shad(g(\zero),\sigma_0)$ for some $\sigma_0 > 0$ independent of $\eta$ and $t$. Then by the Bounded Distortion Lemma \ref{lemmaboundeddistortion}, we have
\begin{align*}
B\left(\bp, e^{-(t - t_\xi)}/(2C)\right) &\subset g^{-1}(B(\xi,e^{-t}/2))\\ 
&\subset g^{-1}(B(\eta,e^{-t}))\\
&\subset g^{-1}(B(\xi,2e^{-t}))\\
&\subset B\left(\bp,2C e^{-(t - t_\xi)}\right)
\end{align*}
for some $C > 0$, and thus
\[
e^{-\delta t_\xi} \mu\left(B\left(\bp, e^{-(t - t_\xi)}/(2C)\right)\right) \lesssim_\times  \mu(B(\eta,e^{-t})) \lesssim_\times e^{-\delta t_\xi} \mu\left(B\left(\bp,2C e^{-(t - t_\xi)}\right)\right).
\]
Combining with \eqref{muBbprbounds2} completes the proof.
\end{proof}

\begin{lemma}
\label{lemmaglobalmeasure2}
For all $\eta\in \Lambda\butnot\{\bp\}$ and $3R_\bp \leq R\leq \Dist_\bp(\zero,\eta)/2$,
\[
\Dist_\bp(\zero,\eta)^{-2\delta} \NN_\bp(R/2) \lesssim_\times \mu(B_\bp(\eta,R)) \lesssim_\times \Dist_\bp(\zero,\eta)^{-2\delta} \NN_\bp(2R).
\]
\end{lemma}
\begin{proof}
Since $\eta\in \Lambda\butnot\{\bp\} \subset G_\bp(S_\bp)$, there exists $h_\eta\in G_\bp$ such that $\eta\in h_\eta(S_\bp)$. Since
\[
\bigcup_{\substack{h\in G_\bp \\ \|h\|_\bp \leq R - R_\bp}} h_\eta h(S_\bp) \subset B_\bp(\eta,R) \subset \bigcup_{\substack{h\in G_\bp \\ \|h\|_\bp \leq R + R_\bp}} h_\eta h(S_\bp),
\]
Lemma \ref{lemmaglobalmeasure1} gives
\[
\sum_{\substack{h\in G_\bp \\ \|h\|_\bp \leq R - R_\bp}} \|h_\eta h\|_\bp^{-2\delta} \lesssim_\times \mu(B_\bp(\eta,R)) \lesssim_\times \sum_{\substack{h\in G_\bp \\ \|h\|_\bp \leq R + R_\bp}} \|h_\eta h\|_\bp^{-2\delta}.
\]
The proof will be complete if we can show that for each $h\in G_\bp$ such that $\|h\|_\bp \leq R + R_\bp$, we have
\begin{equation}
\label{heta}
\|h_\eta h\|_\bp \asymp_\times \Dist_\bp(\zero,\eta).
\end{equation}
And indeed,
\[
\Dist_\bp(\eta,h_\eta h(\zero)) \leq \Dist_\bp(\eta,h_\eta(\zero)) + \|h\|_\bp \leq R_\bp + (R + R_\bp) \leq \frac56 \Dist_\bp(\zero,\eta),
\]
demonstrating \eqref{heta} with an implied constant of $6$.
\end{proof}



\begin{corollary}
\label{corollaryglobalmeasure21}
For all $\eta\in\Lambda\butnot\{\bp\}$ and $6R_\bp \Dist(\bp,\eta)^2 \leq r \leq \Dist(\bp,\eta)/4$, we have
\[
\Dist(\bp,\eta)^{2\delta} \NN_\bp\left(\frac{r}{4\Dist(\bp,\eta)^2}\right)
\lesssim_\times \mu(B(\eta,r))
\lesssim_\times \Dist(\bp,\eta)^{2\delta} \NN_\bp\left(\frac{4r}{\Dist(\bp,\eta)^2}\right).
\]
\end{corollary}


\begin{proof}
By \eqref{GMVT2}, for every $\zeta\in B_\bp\left(\eta,\frac{r}{\Dist(\bp,\eta)(\Dist(\bp,\eta) + r)}\right)$ we have that
\begin{align*}
\Dist(\eta,\zeta) = \frac{\Dist_\bp(\eta,\zeta)}{\Dist_\bp(\zero,\eta)\Dist_\bp(\zero,\zeta)} 
&\leq \frac{\Dist_\bp(\eta,\zeta)}{\Dist_\bp(\zero,\eta)(\Dist_\bp(\zero,\eta) - \Dist_\bp(\eta,\zeta))}\\
&\leq \frac{\frac{r}{\Dist(\bp,\eta)(\Dist(\bp,\eta) + r)}}{\Dist_\bp(\zero,\eta)\left(\Dist_\bp(\zero,\eta) - \frac{r}{\Dist(\bp,\eta)(\Dist(\bp,\eta) + r)}\right)}\\ 
&= r. 
\end{align*}
Analogously, \eqref{GMVT2} also implies that for every $\zeta \in B(\eta,r)$ we have 
\begin{align*}
\Dist_\bp(\eta,\zeta) 
&= \frac{\Dist(\eta,\zeta)}{\Dist(\bp,\eta)\Dist(\bp,\zeta)}\\ 
&\leq \frac{r}{\Dist(\bp,\eta) \left( \Dist(\bp,\eta) - r \right)} \cdot
\end{align*}
Combining these inequalities gives us that 
\[
B_\bp\left(\eta,\frac{r}{\Dist(\bp,\eta)(\Dist(\bp,\eta) + r)}\right)
\subset B(\eta,r)
\subset B_\bp\left(\eta,\frac{r}{\Dist(\bp,\eta)(\Dist(\bp,\eta) - r)}\right)
\]
Now since $r\leq \Dist(\bp,\eta)/4$, we have
\[
B_\bp\left(\eta,\frac{r}{2\Dist(\bp,\eta)^2}\right)
\subset B(\eta,r)
\subset B_\bp\left(\eta,\frac{2r}{\Dist(\bp,\eta)^2}\right).
\]
On the other hand, since $6R_\bp \Dist(\bp,\eta)^2 \leq r \leq \Dist(\bp,\eta)/4$, we have
\[
3R_\bp \leq \frac{r}{2\Dist(\bp,\eta)^2} \leq \frac{2r}{\Dist(\bp,\eta)^2} \leq \frac{\Dist(\bp,\eta)}{2}
\]
whereupon Lemma \ref{lemmaglobalmeasure2} completes the proof.
\end{proof}
%\begin{proof}
%By \eqref{GMVT2}, we have
%\[
%B_\bp\left(\eta,\frac{r}{\Dist(\bp,\eta)(\Dist(\bp,\eta) + r)}\right)
%\subset B(\eta,r)
%\subset B_\bp\left(\eta,\frac{r}{\Dist(\bp,\eta)(\Dist(\bp,\eta) - r)}\right);\Footnote{If $\zeta\in B_\bp\left(\eta,\frac{r}{\Dist(\bp,\eta)(\Dist(\bp,\eta) + r)}\right)$, then 
%\begin{align*}
%\Dist(\eta,\zeta) = \frac{\Dist_\bp(\eta,\zeta)}{\Dist_\bp(\zero,\eta)\Dist_\bp(\zero,\zeta)} 
%&\leq \frac{\Dist_\bp(\eta,\zeta)}{\Dist_\bp(\zero,\eta)(\Dist_\bp(\zero,\eta) - \Dist_\bp(\eta,\zeta))}\\
%&\leq \frac{\frac{r}{\Dist(\bp,\eta)(\Dist(\bp,\eta) + r)}}{\Dist_\bp(\zero,\eta)\left(\Dist_\bp(\zero,\eta) - \frac{r}{\Dist(\bp,\eta)(\Dist(\bp,\eta) + r)}\right)} = r. 
%\end{align*}
%}
%\]
%since $r\leq \Dist(\bp,\eta)/4$, we have
%\[
%B_\bp\left(\eta,\frac{r}{2\Dist(\bp,\eta)^2}\right)
%\subset B(\eta,r)
%\subset B_\bp\left(\eta,\frac{2r}{\Dist(\bp,\eta)^2}\right).
%\]
%On the other hand, since $6R_\bp \Dist(\bp,\eta)^2 \leq r \leq \Dist(\bp,\eta)/4$, we have
%\[
%3R_\bp \leq \frac{r}{2\Dist(\bp,\eta)^2} \leq \frac{2r}{\Dist(\bp,\eta)^2} \leq \frac{\Dist(\bp,\eta)}{2},
%\]
%whereupon Lemma \ref{lemmaglobalmeasure2} completes the proof.
%\end{proof}

\begin{figure}
\begin{center}
\begin{tabular}{ll}
%\begin{tikzpicture}[line cap=round,line join=round,>=triangle 45,x=1.0cm,y=1.0cm]
\begin{tikzpicture}[line cap=round,line join=round,>=triangle 45,scale=.81]
%\clip(-4.3527228989498425,-4.317699333578228) rectangle (4.5670505753079516,4.22893620649748);
\clip(-2.7,-2.7) rectangle (4.3,4.3);
\draw(0.0,0.0) circle (4.0cm);
\draw(2.342806400391098,0.712964401564811) circle (1.551112790229163cm);
\draw (0.0,0.0)-- (4.0,0.0);
\draw(3.212205575454762,0.0) circle (0.3760152197599791cm);
%\draw (0.0,0.0)-- (3.9725001940685334,0.4682330703030985);
\draw [dash pattern=on 3pt off 3pt] (0.0,0.0)-- (3.9725001940685334,0.4682330703030985);
%\draw (0.0,0.0)-- (3.9725001940685334,-0.4682330703030985);
\draw [dash pattern=on 3pt off 3pt] (0.0,0.0)-- (3.9725001940685334,-0.4682330703030985);
\draw [shift={(0.0,0.0)},line width=1.2000000000000002pt]  plot[domain=-0.11732726409329874:0.11732726409329915,variable=\t]({1.0*4.0*cos(\t r)+-0.0*4.0*sin(\t r)},{0.0*4.0*cos(\t r)+1.0*4.0*sin(\t r)});
\begin{scriptsize}
\draw [fill=black] (0.0,0.0) circle (1pt);
\draw[color=black] (-0.2491340591615706,-0.0014223831473892945) node {$\zero$};
\draw [fill=black] (4.0,0.0) circle (1pt);
\draw[color=black] (4.211681114182143,0.0) node {$\eta$};
\draw [fill=black] (3.825718764158789,1.1678509911642618) circle (1pt);
\draw[color=black] (4.057469971268139,1.2) node {$\xi$};
\draw [fill=black] (3.384151714048233,0.0) circle (1pt);
\draw[color=black] (3.2344150960861695,-0.1380274618110051) node {$\eta_t$};
\end{scriptsize}
\end{tikzpicture}

%\begin{tikzpicture}[line cap=round,line join=round,>=triangle 45,x=1.0cm,y=1.0cm]
\begin{tikzpicture}[line cap=round,line join=round,>=triangle 45,scale=.85]
\clip(-4.152066304477812,-0.497280356328431) rectangle (4.307443018392998,5.909099879057296);
\draw (-3.5,0.0)-- (3.5,0.0);
\draw (-3.5,1.5)-- (3.5,1.5);
\draw [shift={(0.40000000000000036,0.0)}] plot[domain=0.3996780691889681:3.141592653589793,variable=\t]({1.0*2.518491612056709*cos(\t r)+-0.0*2.518491612056709*sin(\t r)},{0.0*2.518491612056709*cos(\t r)+1.0*2.518491612056709*sin(\t r)});
\draw (2.4230669786242873,1.5)-- (2.1163456022508336,1.2434075389985795);
\draw (-2.1184916120567086,0.0)-- (-2.12,5.4);
\draw (2.1163456022508336,1.2434075389985795)-- (2.12,5.4);
\draw [dash pattern=on 1pt off 1pt] (-1.5356104253500993,1.6112765998660834)-- (-2.1189952002490946,1.8028361012686287);
\begin{scriptsize}
\draw [fill=black] (2.72,0.98) circle (1pt);
\draw[color=black] (2.9026704848505824,1.0) node {$\zero$};
\draw [fill=black] (-2.1184916120567086,0.0) circle (1pt);
\draw[color=black] (-2.284181946690644,-0.1267908969326299) node {$\eta$};
\draw [fill=black] (2.4230669786242873,1.5) circle (1pt);
\draw[color=black] (2.58,1.68) node {$\eta_{t_\xi}$};
\draw [fill=black] (2.1163456022508336,1.2434075389985795) circle (1pt);
\draw[color=black] (1.8073209579773947,1.092736906911882) node {$g(\zero)$};
%\draw [fill=black] (-0.13843049435662466,5.4459880548125446) circle (1pt);
\draw[color=black] (-0.1,5.662106906126763) node {$\xi=\infty$};
\draw [fill=black] (-1.5356104253500993,1.6112765998660834) circle (1pt);
\draw[color=black] (-1.2,1.65) node {$\eta_t$};
\draw[color=black] (-3.084181946690644,1.867908969326299) node {$H_{\xi}$};
\end{scriptsize}
\end{tikzpicture}
\end{tabular}

\caption[Estimating measures of balls via ``local'' information]{Cusp excursions in the ball model (left) and upper half-space model (right). Since $\xi = g(\bp)\notin B(\eta,e^{-t})$, our estimate of $\mu(B(\eta,e^{-t}))$ is based on the function $\NN_\bp$, which captures ``local'' information about the cusp $\bp$. In the right-hand picture, the measure of $B(\eta,e^{-t})$ can be estimated by considering the measure from the perspective of $g(\zero)$ of a large ball around $\eta$ taken with respect to the $\Dist_\xi$-metametric.}
\label{figurecuspleft}
\end{center}
\end{figure}



\begin{corollary}[Cf. Figure \ref{figurecuspleft}]
\label{corollaryglobalmeasure22}
Fix $\eta\in\Lambda$ and $t > 0$ such that $\eta_t\in H_\xi$ for some $\xi = g(\bp)\in\Lbp$. If
\begin{equation}
\label{globalmeasure22hypothesis}
\lb \xi|\eta\rb_\zero + \tau \leq t \leq 2\lb \xi|\eta\rb_\zero - t_\xi - \tau,
\end{equation}
then
\begin{equation}
\label{globalmeasure22}
\begin{split}
e^{-\delta (2\lb \xi|\eta\rb_\zero - t_\xi)} \NN_\bp(e^{2\lb \xi|\eta\rb_\zero - t_\xi - t - \sigma})
&\lesssim_\times \mu\big(B(\eta,e^{-t})\big)\\ 
&\lesssim_\times e^{-\delta (2\lb \xi|\eta\rb_\zero - t_\xi)} \NN_\bp(e^{2\lb \xi|\eta\rb_\zero - t_\xi - t + \sigma})
\end{split}
\end{equation}
where $\sigma,\tau > 0$ are independent of $\eta$ and $t$.
\end{corollary}
\begin{proof}
Without loss of generality suppose that $g$ satisfies \eqref{top}, and write $\zeta = g^{-1}(\eta)$. Since $t > t_\xi$, \eqref{ShadcontainsB} guarantees that $B(\eta,e^{-t}) \subset \Shad(g(\zero),\sigma_0)$ for some $\sigma_0 > 0$ independent of $\eta$ and $t$. Then by the Bounded Distortion Lemma \ref{lemmaboundeddistortion}, we have
\[
B(\zeta, e^{-(t - t_\xi)}/C) \subset g^{-1}(B(\eta,e^{-t})) \subset B(\zeta,C e^{-(t - t_\xi)})
\]
for some $C > 0$, and thus
\[
e^{-\delta t_\xi} \mu(B(\zeta, e^{-(t - t_\xi)}/C)) \lesssim_\times  \mu(B(\eta,e^{-t})) \lesssim_\times e^{-\delta t_\xi} \mu(B(\zeta,C e^{-(t - t_\xi)})).
\]
If
\begin{equation}
\label{globalmeasure22hypothesisinternal}
6R_\bp \Dist(\bp,\eta)^2 \leq \frac{e^{-(t - t_\xi)}}{C} \leq C e^{-(t - t_\xi)} \leq \frac{\Dist(\bp,\zeta)}{4},
\end{equation}
then Corollary \ref{corollaryglobalmeasure21} guarantees that
%\[
%e^{-\delta t_\xi} \Dist(\bp,\zeta)^{2\delta} \NN_\bp \lesssim_\times  \mu(B(\eta,e^{-t})) \lesssim_\times e^{-\delta t_\xi} \mu(B(\zeta,C e^{-(t - t_\xi)})).
%\]
\begin{align*}
e^{-\delta t_\xi} \Dist(\bp,\zeta)^{2\delta} \NN_\bp\left(\frac{e^{-(t - t_\xi)}}{4C\Dist(\bp,\zeta)^2}\right)
&\lesssim_\times \mu(B(\eta,e^{-t}))\\
&\lesssim_\times e^{-\delta t_\xi}\Dist(\bp,\zeta)^{2\delta} \NN_\bp\left(\frac{4C e^{-(t - t_\xi)}}{\Dist(\bp,\zeta)^2}\right).
\end{align*}
On the other hand, since $\xi,\eta\in\Shad(g(\zero),\sigma_0)$, the Bounded Distortion Lemma \ref{lemmaboundeddistortion} guarantees that
\[
\Dist(\bp,\zeta) \asymp_\times e^{t_\xi} \Dist(\xi,\eta) = e^{-(\lb \xi|\eta\rb_\zero - t_\xi)}.
\]
Denoting the implied constant by $K$, we deduce \eqref{globalmeasure22} with $\sigma = \log(4CK^2)$. The proof is completed upon observing that if $\tau = \log(4CK\vee 6R_\bp CK^2)$, then \eqref{globalmeasure22hypothesis} implies \eqref{globalmeasure22hypothesisinternal}.
%\[
%\lb \xi|\eta\rb_\zero + \tau \leq t \leq 2\lb \xi|\eta\rb_\zero - t_\xi - \tau \;\;\Rightarrow\;\; \eqref{}
%\]
%and thus
%\[
%e^{\delta t_\xi} \Dist(\xi,\eta)^{2\delta} \NN_\bp\left(\frac{e^{-(t - t_\xi)}}{4C C_2^2 \Dist(\xi,\eta)^2}\right)
%\lesssim_\times \mu(B(\eta,e^{-t}))
%\lesssim_\times e^{\delta t_\xi} \Dist(\xi,\eta)^{2\delta} \NN_\bp\left(\frac{4C C_2^2 e^{-(t - t_\xi)}}{\Dist(\xi,\eta)^2}\right).
%\]
%for some $C_2 > 0$.
\end{proof}

\begin{lemma}[Cf. Lemma \ref{lemmasullivanshadow}]
\label{lemmaCCBglobalmeasure}
Fix $\eta\in\Lambda$ and $t > 0$ such that $\eta_t\notin\bigcup(\scrH)$. Then
\[
\mu(B(\eta,e^{-t})) \asymp_\times e^{-\delta t}.
\]
\end{lemma}
\begin{proof}
By \eqref{sigmadef}, there exists $g\in G$ such that $\dist(g(\zero),\eta_t) \asymp_\plus 0$. By \eqref{ShadcontainsB}, we have $B(\eta,e^{-t}) \subset \Shad(g(\zero),\sigma)$ for some $\sigma > 0$ independent of $\eta,t$. It follows that
\[
\mu(B(\eta,e^{-t})) \asymp_\times e^{-\delta t} \mu(g^{-1}(B(\eta,e^{-t}))).
\]
To complete the proof it suffices to show that $\mu(g^{-1}(B(\eta,e^{-t})))$ is bounded from below. By the Bounded Distortion Lemma \ref{lemmaboundeddistortion},
\[
g^{-1}(B(\eta,e^{-t})) \supset B(g^{-1}(\eta),\epsilon)
\]
for some $\epsilon > 0$ independent of $\eta,t$. Now since $G$ is of compact type, we have
\[
\inf_{x\in\Lambda} \mu(B(x,\epsilon)) \geq \min_{x\in S_{\epsilon/2}} \mu(B(x,\epsilon/2)) > 0
\]
where $S_{\epsilon/2}$ is a maximal $\epsilon/2$-separated subset of $\Lambda$. This completes the proof.
\end{proof}

We are now ready to prove Theorem \ref{theoremglobalmeasure}:

\begin{proof}[Proof of Theorem \ref{theoremglobalmeasure}]
Let $\sigma_0$ denote the implied constant of \eqref{betatinterpretation}. Then by \eqref{metat}, for all $\eta\in\Lambda$, $t > 0$, and $\xi\in\Lbp$,
\begin{equation}
\label{metatsimple}
m(\eta,t) = \begin{cases}
e^{-\delta t_\xi} [\II_\bp(e^{t - t_\xi - \theta}) + \mu(\bp)] & t_\xi + \sigma_0 \leq t \leq \lb \xi|\eta\rb_\zero\\
e^{-\delta (2\lb \xi|\eta\rb_\zero - t_\xi)} \NN_\bp(e^{2\lb \xi|\eta\rb_\zero - t - t_\xi - \theta}) & \lb \xi|\eta\rb_\zero < t \leq 2\lb \xi|\eta\rb_\zero - t_\xi - \sigma_0\\
\text{unknown} & \text{otherwise}
\end{cases}
\end{equation}
Applying this formula to Corollaries \ref{corollaryglobalmeasure12} and \ref{corollaryglobalmeasure22} yields the following:

\begin{lemma}
\label{lemmaglobalmeasureinternal}
There exists $\tau\geq \sigma_0$ such that for all $\eta\in\Lambda$ and $t > 0$.
\begin{itemize}
\item[(i)] If for some $\xi$, $t_\xi + \tau \leq t \leq \lb \xi|\eta\rb_\zero - \tau$, then \eqref{globalmeasureformula} holds.
\item[(ii)] If for some $\xi$, $\lb \xi|\eta\rb_\zero + \tau \leq t \leq 2\lb \xi|\eta\rb_\zero - t_\xi - \tau$, then \eqref{globalmeasureformula} holds.
\end{itemize}
\end{lemma}
Now fix $\eta\in\Lambda$, and let
\[
A = \left\{t > 0 : \eta_t\notin\bigcup(\scrH)\right\} \cup \bigcup_{\xi\in\Lbp} [t_\xi + \tau,\lb \xi|\eta\rb_\zero - \tau]\cup \bigcup_{\xi\in\Lbp}  [\lb \xi|\eta\rb_\zero + \tau,2\lb \xi|\eta\rb_\zero - t_\xi - \tau].
\]
Then by Lemmas \ref{lemmaCCBglobalmeasure} and \ref{lemmaglobalmeasureinternal}, $\eqref{globalmeasureformula}_{\sigma = \tau}$ holds for all $t\in A$.
\begin{claim}
Every interval of length $2\tau$ intersects $A$.
\end{claim}
\begin{subproof}
If $[s - \tau,s + \tau]$ does not intersect $A$, then by connectedness, there exists $\xi\in\Lbp$ such that $\eta_t\in H_\xi$ for all $t\in [s - \tau, s + \tau]$. By Lemma \ref{lemmabetatinterpretation}, the fact that $\eta_{s\pm\tau}\in H_\xi$ implies that $t_\xi \leq s \leq 2\lb \xi|\eta\rb_\zero - t_\xi$ (since $\tau\geq \sigma_0$). If $s \leq \lb \xi|\eta\rb_\zero$, then $[s - \tau,s + \tau]\cap [t_\xi + \tau,\lb \xi|\eta\rb_\zero - \tau] \neq \emptyset$, while if $s \geq \lb \xi|\eta\rb_\zero$, then $[s - \tau,s + \tau]\cap [\lb \xi|\eta\rb_\zero + \tau,2\lb \xi|\eta\rb_\zero - t_\xi - \tau] \neq \emptyset$.
\end{subproof}
Thus for all $t > 0$, there exist $t_\pm\in A$ such that $t - 2\tau \leq t_- \leq t \leq t_+ \leq t - 2\tau$; then
\begin{align*}
m(\eta,t + 3\tau)
\lesssim_\times m(\eta,t_+ + \tau)
&\lesssim_\times \mu\big(B(\eta,e^{-t_+})\big)\\
&\leq_\pt \mu\big(B(\eta,e^{-t})\big)\\
&\leq_\pt \mu\big(B(\eta,e^{-t_-})\big)
\lesssim_\times m(\eta,t_- - \tau)
\lesssim_\times m(\eta,t - 3\tau),
\end{align*}
i.e. $\eqref{globalmeasureformula}_{\sigma = 3\tau}$ holds.
%Now let $\eta\in\Lambda$ and $t > 0$ be arbitrary. If $\eta_t\notin\bigcup(\scrH)$, then combining Lemma \ref{lemmaCCBglobalmeasure} with Proposition \ref{propositionmetadecreasing} demonstrates \eqref{globalmeasureformula}. So suppose $\eta_t\in H_\xi$ for some $\xi\in\Lbp$. If $\lb \xi|\eta\rb_\zero - \sigma < t < \lb \xi|\eta\rb_\zero + \sigma$, then
%\[
%m(\eta,\lb \xi|\eta\rb_\zero + 2\sigma)
%\lesssim_\times \mu\big(B(\eta,e^{-(\lb \xi|\eta\rb_\zero + \sigma)})\big)
%\leq \mu\big(B(\eta,e^{-t})\big)
%\leq \mu\big(B(\eta,e^{-(\lb \xi|\eta\rb_\zero - \sigma)})\big)
%\lesssim_\times m(\eta,\lb \xi|\eta\rb_\zero - 2\sigma),
%\]
%demonstrating $\eqref{globalmeasureformula}_{\sigma = 3\sigma}$. If $t\leq t_\xi + \sigma$, then
%\[
%m(\eta,t_\xi + 2\sigma)
%\lesssim_\times \mu\big(B(\eta,e^{-(t_\xi + \sigma)})\big)
%\leq \mu\big(B(\eta,e^{-t})\big)
%\leq \mu\big(B(\eta,e^{-(\lb \xi|\eta\rb_\zero - \sigma)})\big)
%\lesssim_\times m(\eta,\lb \xi|\eta\rb_\zero - 2\sigma),
%\]
%Let $\sigma,\tau$ be large enough so that Corollaries \ref{corollaryglobalmeasure12} and \ref{corollaryglobalmeasure22} hold. Fix $\eta\in\Lambda$ and $t > 0$. If $\eta_t\notin\bigcup(\scrH)$, then $\mu(B(\eta,e^{-t})) \asymp_\times m(\eta,t)$ by Lemma \ref{lemmaCCBglobalmeasure}, and Proposition \ref{propositionmetadecreasing} finishes the proof. So suppose $\eta_t\in H_\xi$ for some $\xi = g(\bp)\in\Lbp$.
%\begin{itemize}
%\item[Case 1:] $t\leq \lb \xi|\eta\rb_\zero - \max(\log(2),\theta + \sigma)$. In this case, let $t_\pm = t \pm \sigma + \theta$; by Corollary \ref{corollaryglobalmeasure12} we have
%\begin{equation}
%\label{globalmeasure12T}
%e^{-\delta t_\xi} [\II_\bp(e^{t_+ - t_\xi - \theta}) + \mu(\bp)] \lesssim_\times \mu\big(B(\eta,e^{-t})\big) \lesssim_\times e^{-\delta t_\xi} [\II_\bp(e^{t_- - t_\xi - \theta}) + \mu(\bp)].
%\end{equation}
%On the other hand, since $t_\pm \leq \lb \xi|\eta\rb_\zero$, by \eqref{metat} and Lemma \ref{lemmabetatinterpretation} we have
%\[
%m(\eta,t_\pm) = e^{-\delta t_\xi} [\II_\bp(e^{t_\pm - t_\xi - \theta}) + \mu(\bp)] \text{ unless } t_\pm \leq t_\xi + \sigma_0,
%\]
%where $\sigma_0$ is the implied constant of \eqref{betatinterpretation}. But if $t_\pm \leq t_\xi + \sigma_0$, then
%\[
%e^{\delta t_\xi} m(\eta,t_\pm) \asymp_\times \II_\bp(e^{t_\pm - t_\xi - \theta}) + \mu(\bp)
%\]
%since both sides are bounded from above and below. Thus in any case we have
%\[
%m(\eta,t_\pm) \asymp_\times e^{-\delta t_\xi} [\II_\bp(e^{t_\pm - t_\xi - \theta}) + \mu(\bp)];
%\]
%plugging into \eqref{globalmeasure12T} and combining with Proposition \ref{propositionmetadecreasing} completes the proof.
%\item[Case 2:] $t > \lb \xi|\eta\rb_\zero + \max(\tau,\sigma - \theta)$. In this case, let $t_\pm = t \pm \sigma + \theta$; by Corollary \ref{corollaryglobalmeasure22} we have
%\begin{equation}
%\label{globalmeasure22T}
%e^{-\delta (2\lb \xi|\eta\rb_\zero - t_\xi)} \NN_\bp(e^{2\lb \xi|\eta\rb_\zero - t_\xi - t_+ + \theta})
%\lesssim_\times \mu\big(B(\eta,e^{-t})\big)
%\lesssim_\times e^{-\delta (2\lb \xi|\eta\rb_\zero - t_\xi)} \NN_\bp(e^{2\lb \xi|\eta\rb_\zero - t_\xi - t_- - \theta}) \text{ unless } t_- 
%\end{equation}
%On the other hand, since $t_\pm > \lb \xi|\eta\rb_\zero$, by \eqref{metat} and Lemma \ref{lemmabetatinterpretation}
%
%\item[Case 3:] $\lb \xi|\eta\rb_\zero - \max(\log(2),\theta + \sigma) < t \leq \lb \xi|\eta\rb_\zero + \max(\tau,\sigma - \theta)$.
%
%\end{itemize}
\end{proof}




%Suppose $\eta_t\in H_\xi$ for some $\xi = g(\bp)\in\Lbp$, $\bp\in P$, $g\in G$ satisfying \eqref{top}. Note that since $\eta_t\in H_\xi$, we have $t_\xi < \busemann_\xi(\zero,\eta_t) \leq t$. Let $\zeta = g^{-1}(\eta)$. Fix $\sigma > 0$ large enough so that $B(\eta,3 e^{-t}) \subset \Shad(g(\zero),\sigma)$ for all $\eta\in\del X$ and $t > 0$ (cf. \eqref{ShadcontainsB}, \eqref{top}\internal). Then by the Bounded Distortion Lemma \ref{lemmaboundeddistortion},\begin{equation}
%\label{globalmeasureproof1}
%\mu(B(\eta,e^{-t})) \asymp_\times e^{-\delta\dogo g} \mu(g^{-1}(B(\eta,e^{-t}))) \asymp_\times e^{-\delta t_\xi} \mu(g^{-1}(B(\eta,e^{-t}))).
%\end{equation}
%The Bounded Distortion Lemma also gives
%\begin{equation}
%\label{globalmeasureproof2}
%B(\zeta,e^{-(t - t_\xi)}/C) \subset g^{-1}(B(\eta,e^{-t})) \subset B(\zeta,C e^{-(t - t_\xi)})
%\end{equation}
%and
%\begin{equation}
%\label{globalmeasureproof3}
%\frac1C \leq \frac{\Dist(g(\bp),\eta)}{e^{-t_\xi}\Dist(\bp,\zeta)} \leq C
%\end{equation}
%for some constant $C \geq 1$.
%
%If $\Dist(g(\bp),\eta) \leq e^{-t}/2$, then
%\begin{align*}
%B(\bp,e^{-(t - t_\xi)}/(2C))
%&\subset g^{-1}(B(g(\bp),e^{-t}/2))\\
%&\subset g^{-1}(B(\eta,e^{-t}))\\
%&\subset g^{-1}(B(g(\bp),2 e^{-t}))
%\subset B(\bp,2C e^{-(t - t_\xi)}),
%\end{align*}
%and thus by Corollary \ref{corollaryglobalmeasure1},
%\[
%\II_\bp(4C e^{t - t_\xi})
%\lesssim_\times \mu(g^{-1}(B(\eta,e^{-t})))
%\lesssim_\times \II_\bp(e^{t - t_\xi}/(4C)).
%\]
%Combining with \eqref{betatasymp} completes the proof in this case.
%
%If
%\begin{equation}
%\label{farend}
%4C e^{-(t - t_\xi)} \leq \Dist(\bp,\zeta) \leq \sqrt{e^{-(t - t_\xi)}/(6C R_\bp)},
%\end{equation}
%then by Corollary \ref{corollaryglobalmeasure2}, we have
%\begin{align*}
%\Dist(\bp,\zeta)^{2\delta} \NN_\bp\left(\frac{e^{-(t - t_\xi)}}{4C\Dist(\bp,\zeta)^2}\right)
%&\lesssim_\times \mu(B(\zeta,e^{-(t - t_\xi)}/C))\\
%&\lesssim_\times \mu(g^{-1}(B(\eta,e^{-t})))\\
%&\lesssim_\times \mu(B(\zeta,C e^{-(t - t_\xi)}))\\
%&\lesssim_\times \Dist(\bp,\zeta)^{2\delta} \NN_\bp\left(\frac{4C e^{-(t - t_\xi)}}{\Dist(\bp,\zeta)^2}\right).
%\end{align*}
%Combining with \eqref{betatasymp} completes the proof in this case.

\section{Groups for which $\mu$ is doubling}
\label{subsectiondoubling}

Recall that a measure $\mu$ is said to be \emph{doubling} if for all $\eta\in\Supp(\mu)$ and $r > 0$, $\mu(B(\eta,2r)) \asymp_\times \mu(B(\eta,r))$. In the Standard Case, the Global Measure Formula implies that the $\delta$-conformal measure of a geometrically finite group is always doubling (Example \ref{exampledoubling}). However, in general there are geometrically finite groups whose $\delta$-conformal measures are not doubling (Example \ref{exampledoublingnondoubling}). It is therefore of interest to determine necessary and sufficient conditions on a geometrically finite group for its $\delta$-conformal measure to be doubling. The Global Measure Formula immediately yields the following criterion:% Abusing terminology, we will call a geometrically finite group $G$ ``doubling'' if and only if its $\delta$-conformal measure is doubling.


\begin{lemma}
$\mu$ is doubling if and only if the function $m$ satisfies
\begin{equation}
\label{mdoubling}
m(\eta,t + \tau) \asymp_{\times,\tau} m(\eta,t - \tau) \all \eta\in\Lambda \all t,\tau > 0.
\end{equation}
\end{lemma}
\begin{proof}
If \eqref{mdoubling} holds, then \eqref{globalmeasureformula} reduces to
\begin{equation}
\label{globalmeasureformula2}
\mu(B(\eta,e^{-t})) \asymp_\times m(\eta,t),
\end{equation}
and then \eqref{mdoubling} shows that $\mu$ is doubling. On the other hand, if $\mu$ is doubling, then \eqref{globalmeasureformula} implies that
\[
m(\eta,t - \tau) \lesssim_\times \mu(B(\eta,e^{-(t - \tau - \sigma)})) \asymp_\times \mu(B(\eta,e^{-(t + \tau + \sigma)})) \lesssim_\times m(\eta,t + \tau);
\]
combining with Proposition \ref{propositionmetadecreasing} shows that \eqref{mdoubling} holds.
\end{proof}

Of course, the criterion \eqref{mdoubling} is not very useful by itself, since it refers to the complicated function $m$. In what follows we find more elementary necessary and sufficient conditions for doubling. First we must introduce some terminology.

\begin{definition}
\label{definitiondoubling}
A function $f:\CO 1\infty\to\CO 1\infty$ is called \emph{doubling} if there exists $\beta > 1$ such that
\begin{equation}
\label{dexp1}
f(\beta R) \lesssim_{\times,\beta} f(R) \all R\geq 1,
\end{equation}
and \emph{codoubling} if there exists $\beta > 1$ such that
\begin{equation}
\label{dexp2}
f(\beta R) - f(R) \gtrsim_{\times,\beta} f(R) \all R\geq 1.
\end{equation}
\end{definition}

\begin{observation}
\label{observationcodoubling}
If there exists $\beta > 1$ such that
\[
\NN_\bp(\beta R) > \NN_\bp(R) \all R\geq 1,
\]
then $\NN_\bp$ is codoubling.
\end{observation}
\begin{proof}
Fix $R\geq 1$; there exists $h\in G_\bp$ such that $2R < \|h\|_\bp \leq 2\beta R$. We have
\[
h\{j\in G_\bp : \|j\|_\bp \leq R\} \subset \{j\in G_\bp : R < \|j\|_\bp \leq (2\beta + 1) R\},
\]
and taking cardinalities gives
\[
\NN_\bp(R) \leq \NN_\bp\big((2\beta + 1) R\big) - \NN_\bp(R).
\]
\end{proof}

We are now ready to state a more elementary characterization of when $\mu$ is doubling:

%\begin{definition}
%\label{definitionregular}
%A parabolic group $H$ with fixed point $\bp$ is called \emph{regular} if there exists $\lambda > 1$ such that for all $R\geq 1$,
%\[
%\NN_\bp(\lambda R) > \NN_\bp(R).
%\]
%\end{definition}
%
%\begin{observation}
%If $H$ has an element of infinite order, then $H$ is regular.
%\end{observation}
%\begin{proof}
%If $h\in H$ has infinite order, let $n$ be minimal such that $\|h^n\|_\bp > R$; then $\|h^n\|\leq 2R$.
%\end{proof}

\begin{proposition}
\label{propositiondoubling}
$\mu$ is doubling if and only if all of the following hold:
\begin{itemize}
\item[(I)] For all $\bp\in P$, $\NN_\bp$ is both doubling and codoubling.
\item[(II)] For all $\bp\in P$ and $R\geq 1$,
\begin{equation}
\label{INasymp2}
\II_\bp(R) \asymp_\times R^{-2\delta} \NN_\bp(R).
\end{equation}
\item[(III)] $G$ is of divergence type.
\end{itemize}
Moreover, \text{(II)} can be replaced by
\begin{itemize}
\item[(II$'$)] For all $\bp\in P$ and $R\geq 1$,
\begin{equation}
\label{INasymp3}
%\II_\bp(R) \lesssim_\times R^{-2\delta} \NN_\bp(R)
\w\II_\bp(R) := \sum_{k = 0}^\infty e^{-2\delta k} \NN_\bp(e^k R) \asymp_\times \NN_\bp(R).
\end{equation}
\end{itemize}
\end{proposition}

\begin{proof}[Proof that \text{(I)-(III)} imply $\mu$ doubling]
Fix $\eta\in\Lambda$ and $t,\tau > 0$, and we will demonstrate \eqref{mdoubling}. By (II), (III), and Observation \ref{observationnonatomic}, we have
\begin{equation}
\label{GMFdoubling1}
\begin{split}
m(\eta,t) &\asymp_\times \begin{cases}
e^{-\delta t} & \eta_t\notin\bigcup(\scrH) \\
e^{-\delta t_\xi} e^{-2\delta(t - t_\xi - \theta)}\NN_\bp(e^{t - t_\xi - \theta}) & \eta_t\in H_\xi \text{ and }t \leq \lb \xi|\eta\rb_\zero\\
e^{-\delta (2\lb \xi|\eta\rb_\zero - t_\xi)} \NN_\bp(e^{2\lb \xi|\eta\rb_\zero - t - t_\xi - \theta}) & \eta_t\in H_\xi \text{ and } t > \lb \xi|\eta\rb_\zero
\end{cases}\\
&\asymp_\times e^{-\delta t} \begin{cases}
1 & \eta_t\notin\bigcup(\scrH)\\
e^{-\delta b(\eta,t)} \NN_\bp(e^{b(\eta,t) - \theta}) & \eta_t \in H_{g(\bp)}
\end{cases}
\end{split}
\end{equation}
where $b(\eta,t)$ is as in \eqref{betat}. Let $t_\pm = t \pm \tau$. We split into two cases:
\begin{itemize}
\item[Case 1:] $\eta_{t_+},\eta_{t_-}\in H_{g(\bp)}$ for some $g(\bp)\in \Lbp$. In this case, \eqref{mdoubling} follows from Corollary \ref{corollarybetat} together with the fact that $\NN_\bp$ is doubling.
\item[Case 2:] $\eta_{t + s}\notin \bigcup(\scrH)$ for some $s\in [-\tau,\tau]$. In this case, Corollary \ref{corollarybetat} shows that $b(\eta,t_\pm) \asymp_{\plus,\tau} 0$ and thus
\[
m(\eta,t_+) \asymp_{\times,\tau} e^{-\delta t} \asymp_{\times,\tau} m(\eta,t_-).
\]
\end{itemize}
\end{proof}
\noindent Before continuing the proof of Proposition \ref{propositiondoubling}, we observe that
\begin{equation*}
%\label{wIasymp}
\begin{split}
\II_\bp(R) + R^{-2\delta}\NN_\bp(R) \asymp_\times \sum_{h\in G_\bp} (R\vee \|h\|_\bp)^{-2\delta}
&\asymp_\times \sum_{h\in G_\bp} \sum_{k = 1}^\infty (e^k R)^{-2\delta} [e^k R \geq \|h\|_\bp]\\
&=_\pt \sum_{k = 1}^\infty (e^k R)^{-2\delta} \NN_\bp(e^k R)\\
&= R^{-2\delta} \w\II_\bp(R).
\end{split}
\end{equation*}
In particular, it follows that \eqref{INasymp3} is equivalent to
\begin{equation}
\label{INasymp4}
\II_\bp(R) \lesssim_\times R^{-2\delta} \NN_\bp(R).
\end{equation}
\begin{proof}[Proof that \text{(I) and (II$'$)} imply \text{(II)}]
Since $\NN_\bp$ is codoubling, let $\beta > 1$ be as in \eqref{dexp2}. Then
\[
\II_\bp(R) \geq \sum_{\substack{h\in G_\bp \\ R < \|h\|_\bp \leq \beta R}} (\beta R)^{-2\delta} = (\beta R)^{-2\delta} (\NN_\bp(\beta R) - \NN_\bp(R)) \gtrsim_{\times,\beta} R^{-2\delta} \NN_\bp(R).
\]
Combining with \eqref{INasymp4} completes the proof.
\end{proof}
\begin{proof}[Proof that $\mu$ doubling implies \text{(I)-(III)} and \text{(II$'$)}]
Since a doubling measure whose topological support is a perfect set cannot have an atomic part, we must have $\mu(P) = 0$ and thus by Observation \ref{observationnonatomic}, (III) holds. Since
\[
m(\bp,t) \asymp_{\times,\bp} \II_\bp(e^{t - t_0 - \theta}) + \mu(\bp) = \II_\bp(e^{t - t_0 - \theta})
\]
for all sufficiently large $t$, setting $\eta = \bp$ in \eqref{mdoubling} shows that the function $\II_\bp$ is doubling.

Fix $\eta\in\Lambda\butnot\{\bp\}$. Let $\sigma_0 > 0$ denote the implied constant of \eqref{betatinterpretation}. For $s\in [t_0 + \sigma_0 + \tau, \lb \bp|\eta\rb_\zero - \tau]$, plugging $t = 2\lb \bp|\eta\rb_\zero - s$ into \eqref{mdoubling} and simplifying using $\eqref{metatsimple}_{\xi = \bp}$ shows that
\begin{equation}
\label{Nbps}
\NN_\bp(e^{s - \tau - t_0 - \theta}) \asymp_{\times,\tau} \NN_\bp(e^{s + \tau - t_0 - \theta}).
\end{equation}
Since $\lb \bp|\eta\rb_\zero$ can be made arbitrarily large, \eqref{Nbps} holds for all $s\geq t_0 + \sigma_0 + \tau$. It follows that $\NN_\bp$ is doubling.

Next, we compare the values of $m(\eta,\lb \bp|\eta\rb_\zero \pm\tau)$. This gives (assuming $\lb \bp|\eta\rb_\zero > t_0 + \sigma_0 + \tau$)
\[
e^{-\delta t_0} \II_\bp(e^{\lb \bp|\eta\rb_\zero - \tau - t_0 - \theta}) \asymp_\times e^{-\delta (2\lb \bp|\eta\rb_\zero - t_0)} \NN_\bp(e^{\lb \bp|\eta\rb_\zero - \tau - t_0 - \theta}).
\]
Letting $R_\eta = \exp(\lb \bp|\eta\rb_\zero - \tau - t_0 - \theta)$, we have
\begin{equation}
\label{INasymp6}
\II_\bp(R_\eta) \asymp_\times R_\eta^{-2\delta} \NN_\bp(R_\eta).
\end{equation}
Now fix $\zeta\in\Lambda\butnot\{\bp\}$ and $h\in G_\bp$, and let $\eta = h(\zeta)$. Then $\Dist_\bp(h(\zero),\eta) \asymp_{\plus,\zeta} 0$, and thus the triangle inequality gives
\[
1 \leq \Dist_\bp(\zero,\eta) \asymp_{\plus,\zeta} \|h\|_\bp \geq 1,
\]
and so $R_\eta \asymp_\times \Dist_\bp(\zero,\eta) \asymp_{\times,\zeta} \|h\|_\bp$. Combining with \eqref{INasymp6} and the fact that the functions $\II_\bp$ and $\NN_\bp$ are doubling, we have
\begin{equation}
\label{INasymp7}
\II_\bp(\|h\|_\bp) \asymp_\times \|h\|_\bp^{-2\delta} \NN_\bp(\|h\|_\bp)
\end{equation}
for all $h\in G_\bp$.

Now fix $1\leq R_1 < R_2$ such that $\|h_i\|_\bp = R_i$ for some $h_1,h_2\in G_\bp$, but such that the formula $R_1 < \|h\|_\bp < R_2$ is not satisfied for any $h\in G_\bp$. Then
\[
\lim_{R\searrow R_1} \II_\bp(R) = \lim_{R\nearrow R_2} \II_\bp(R) \text{ and } \lim_{R\searrow R_1} \NN_\bp(R) = \lim_{R\nearrow R_2} \NN_\bp(R).
%\II_\bp(2R_1) = \II_\bp(R_2/2) \text{ and } \NN_\bp(2R_1) = \NN_\bp(R_2/2).
\]
On the other hand, applying \eqref{INasymp7} with $h = h_1,h_2$ gives
\[
\II_\bp(R_i) \asymp_\times R_i^{-2\delta} \NN_\bp(R_i).
\]
Since $\II_\bp$ and $\NN_\bp$ are doubling, we have
\begin{align*}
%R_1^{-2\delta}
%\asymp_\times \frac{\II_\bp(R_1)}{\NN_\bp(R_1)}
%\asymp_\times \frac{\II_\bp(2R_1)}{\NN_\bp(2R_1)}
%= \frac{\II_\bp(R_2/2)}{\NN_\bp(R_2/2)}
%\asymp_\times \frac{\II_\bp(R_2)}{\NN_\bp(R_2)}
%\asymp_\times R_2^{-2\delta}
R_1^{-2\delta}
\asymp_\times \frac{\II_\bp(R_1)}{\NN_\bp(R_1)}
\asymp_\times \frac{\lim_{R\searrow R_1} \II_\bp(R)}{\lim_{R\searrow R_1} \NN_\bp(R)}
&= \frac{\lim_{R\nearrow R_2} \II_\bp(R)}{\lim_{R\nearrow R_2} \NN_\bp(R)}\\
&\asymp_\times \frac{\II_\bp(R_2)}{\NN_\bp(R_2)}
\asymp_\times R_2^{-2\delta}
\end{align*}
and thus $R_1 \asymp_\times R_2$. Since $R_1,R_2$ were arbitrary, Observation \ref{observationcodoubling} shows that $\NN_\bp$ is codoubling. This completes the proof of (I).

It remains to demonstrate (II) and (II$'$). Given any $R\geq 1$, since $\NN_\bp$ is codoubling, we may find $h\in G_\bp$ such that $\|h\|_\bp \asymp_\times R$; combining with \eqref{INasymp7} and the fact that $\II_\bp$ and $\NN_\bp$ are doubling gives \eqref{INasymp2} and \eqref{INasymp4}, demonstrating (II) and (II$'$).
\end{proof}

We note that the proof actually shows the following (cf. \eqref{GMFdoubling1}):

\begin{corollary}
\label{corollaryGMFdoubling}
If $\mu$ is doubling, then
\[
\mu(B(\eta,e^{-t})) \asymp_\times e^{-\delta t} \begin{cases}
1 & \eta_t\notin \bigcup(\scrH)\\
e^{-\delta b(\eta,t)} \NN_\bp(e^{b(\eta,t)}) & \eta_t\in H_{g(\bp)}
\end{cases}
\]
for all $\eta\in\Lambda$, $t > 0$. Here $b(\eta,t)$ is as in \eqref{betat}.
\end{corollary}

Although Proposition \ref{propositiondoubling} is the best necessary and sufficient condition we can give for doubling, in what follows we give necessary conditions and sufficient conditions which are more elementary (Proposition \ref{propositiondoubling2}), although the necessary conditions are not the same as the sufficient conditions. In practice these conditions are usually powerful enough to determine whether any given measure is doubling.

To state the result, we need the concept of the \emph{polynomial growth rate} of a function:

\begin{definition}[Cf. \eqref{growthrate}]
\label{definitiongrowthrate}
The \emph{(polynomial) growth rate} of a function $f:\CO 1\infty \to \CO 1\infty$ is the limit
\[
\dexp(f) := \lim_{\lambda,R\to\infty}\frac{\log f(\lambda R) - \log f(R)}{\log(\lambda)}
\]
if it exists. If the limit does not exist, then the numbers
\begin{align*}
\uexp(f) &:= \limsup_{\lambda,R\to\infty}\frac{\log f(\lambda R) - \log f(R)}{\log(\lambda)}\\
\lexp(f) &:= \liminf_{\lambda,R\to\infty}\frac{\log f(\lambda R) - \log f(R)}{\log(\lambda)}
\end{align*}
are the \emph{upper} and \emph{lower polynomial growth rates} of $f$, respectively.
\end{definition}

%\begin{remark}
%In general, $(\lexp(f),\uexp(f))$ can be any pair of numbers. However, it is not too unreasonable to look for functions for which $\dexp(f)$ exists; 
%
%\end{remark}

\begin{lemma}
\label{lemmadexp}
Let $f:\CO 1\infty\to\CO 1\infty$.
\begin{itemize}
\item[(i)] $f$ is doubling if and only if $\uexp(f) < \infty$.
\item[(ii)] $f$ is codoubling if and only if $\lexp(f) > 0$.
\item[(iii)]
\[
\lexp(f) \leq \liminf_{\lambda\to\infty} \frac{\log f(\lambda)}{\log(\lambda)} \leq \limsup_{\lambda\to\infty} \frac{\log f(\lambda)}{\log(\lambda)} \leq \uexp(f).
\]
In particular, $\lexp(\NN_\bp) \leq 2\delta_\bp \leq \uexp(\NN_\bp)$.
\end{itemize}
\end{lemma}
\begin{proof}[Proof of \text{(i)}]
Suppose that $f$ is doubling, and let $C > 1$ denote the implied constant of \eqref{dexp1}. Iterating gives
\[
f(\beta^n R) \leq C^n f(R) \all n\in\N \all R\geq 1
\]
and thus
\[
f(\lambda R) \lesssim_\times \lambda^{\log_\beta(C)} f(R) \all \lambda > 1 \all R\geq 1.
\]
It follows that $\uexp(f) \leq \log_\beta(C) < \infty$. The converse direction is trivial.
\end{proof}
\begin{proof}[Proof of \text{(ii)}]
Suppose that $f$ is codoubling, and let $C > 1$ denote the implied constant of \eqref{dexp2}. Then
\[
f(\beta R) \geq (1 + C^{-1}) f(R) \all R\geq 1.
\]
Iterating gives
\[
f(\beta^n R) \geq (1 + C^{-1})^n f(R) \all n\in\N \all R\geq 1
\]
and thus
\[
f(\lambda R) \gtrsim_\times \lambda^{\log_\beta(1 + C^{-1})} f(R) \all \lambda > 1 \all R\geq 1.
\]
It follows that $\lexp(f) \geq \log_\beta(1 + C^{-1}) > 0$. The converse direction is trivial.
\end{proof}
\begin{proof}[Proof of \text{(iii)}]
Let $R_n\to\infty$. For each $n\in\N$,
\[
\limsup_{\lambda\to\infty} \frac{\log f(\lambda R_n) - \log f(R_n)}{\log(\lambda)} = \overline s := \limsup_{\lambda\to\infty} \frac{\log f(\lambda)}{\log(\lambda)}\cdot
\]
Thus given $s < \overline s$, we may find a large number $\lambda_n > 1$ such that 
\[
\frac{\log f(\lambda_n R_n) - \log f(R_n)}{\log(\lambda_n)}  > s.
\] 
Since $\lambda_n,R_n\to\infty$ as $n\to\infty$, it follows that $\uexp(f) \geq s$; since $s$ was arbitrary, $\uexp(f) \geq \overline s$. A similar argument shows that $\lexp(f) \leq \underline s$.

Finally, when $f = \NN_\bp$, the equality $\overline s = \underline s = 2\delta_\bp$ is a consequence of \eqref{poincarealternate} and Observation \ref{observationeuclideanparabolicasymp}.
\end{proof}

We can now state our final result regarding criteria for doubling:

\begin{proposition}
\label{propositiondoubling2}
In the following list, \text{(A) \implies (B) \implies (C):}
\begin{itemize}
\item[(A)] For all $\bp\in P$, $0 < \lexp(\NN_\bp) \leq \uexp(\NN_\bp) < 2\delta$.
\item[(B)] $\mu$ is doubling.
\item[(C)] For all $\bp\in P$, $0 < \lexp(\NN_\bp) \leq \uexp(\NN_\bp) \leq 2\delta$.
\end{itemize}
\end{proposition}
\begin{proof}[Proof of \text{(A) \implies (B)}]
Suppose that (A) holds. Then by Lemma \ref{lemmadexp}, (I) of Proposition \ref{propositiondoubling} holds. Since $\delta_\bp \leq \uexp(\NN_\bp)/2 < \delta$ for all $\bp\in P$, Theorem \ref{theoremGFdivergencetype} implies that (III) of Proposition \ref{propositiondoubling} holds. To complete the proof, we need to show that (II$'$) of Proposition \ref{propositiondoubling} holds. Fix $s\in (\uexp(\NN_\bp),2\delta)$. Since $s > \uexp(\NN_\bp)$, we have
\[
\NN_\bp(\lambda R) \lesssim_{\times,s} \lambda^s \NN_\bp(R) \all \lambda > 1,\; R\geq 1
\]
and thus
\[
\NN_\bp(R) \leq \w\II_\bp(R) \lesssim_\times \sum_{k = 0}^\infty e^{-2\delta k} e^{sk} \NN_\bp(R) \asymp_\times \NN_\bp(R),
\]
demonstrating \eqref{INasymp3} and completing the proof.
\end{proof}
\begin{proof}[Proof of \text{(B) \implies (C)}]
Suppose $\mu$ is doubling. By (I) of Proposition \ref{propositiondoubling}, $\lexp(\NN_\bp) > 0$. On the other hand, by \eqref{INasymp3} we have
\[
\lambda^{-2\delta} \NN_\bp(\lambda R) \lesssim_\times \NN_\bp(R) \all \lambda > 1,\; R\geq 1
\]
and thus $\uexp(\NN_\bp) \leq 2\delta$.
\end{proof}

Proposition \ref{propositiondoubling} shows that if $G$ is a geometrically finite group with $\delta$-conformal measure $\mu$, then the question of whether $\mu$ is doubling is determined entirely by its parabolic subgroups $(G_\bp)_{\bp\in P}$ and its Poincar\'e set $\Delta_G$. A natural question is when the second input can be removed, that is: if we are told what the parabolic subgroups $(G_\bp)_{\bp\in P}$ are, can we sometimes determine whether $\mu$ is doubling without looking at $\Delta_G$? A trivial example is that if $\lexp(\NN_\bp) = 0$ or $\uexp(\NN_\bp) = \infty$ for some $\bp\in P$, then we automatically know that $\mu$ is not doubling. Conversely, the following definition and proposition describe when we can deduce that $\mu$ is doubling:

\begin{definition}
\label{definitionpredoubling}
A parabolic group $H\leq\Isom(X)$ with global fixed point $\bp\in\del X$ is \emph{pre-doubling} if
\begin{equation}
\label{predoubling}
0 < \lexp(\NN_{\EE_\bp,H}) \leq \uexp(\NN_{\EE_\bp,H}) = 2\delta_H < \infty
\end{equation}
and $H$ is of divergence type.
\end{definition}

%\begin{observation}
%\label{observationpredoubling}
%If $G_\bp$ is pre-doubling for every $\bp\in P$, then $\mu$ is doubling.
%\end{observation}
%\begin{proof}
%For all $\bp\in P$, the fact that $G_\bp$ is of divergence type implies that $\delta > \delta_\bp$ (Proposition \ref{propositiondeltagtrdeltap}); combining with \eqref{predoubling} gives $0 < \lexp(\NN_\bp) \leq \uexp(\NN_\bp) < 2\delta$. Proposition \ref{propositiondoubling2} completes the proof.
%\end{proof}

\begin{proposition}
\label{propositionpredoubling}
~
\begin{itemize}
\item[(i)] If $G_\bp$ is pre-doubling for every $\bp\in P$, then $\mu$ is doubling.
\item[(ii)] Let $H\leq\Isom(X)$ be a parabolic subgroup, and let $g\in\Isom(X)$ be a loxodromic isometry such that $\lb g,H\rb$ is a strongly separated Schottky product. Then the following are equivalent:
\begin{itemize}
\item[(A)] $H$ is pre-doubling.
\item[(B)] For every $n\in\N$, the $\delta_n$-quasiconformal measure $\mu_n$ of $G_n = \lb g^n,H\rb$ is doubling. Here we assume that $\delta_n := \delta(G_n) < \infty$.
\end{itemize}
\end{itemize}
\end{proposition}
\begin{proof}[Proof of \text{(i)}]
For all $\bp\in P$, the fact that $G_\bp$ is of divergence type implies that $\delta > \delta_\bp$ (Proposition \ref{propositiondeltagtrdeltap}); combining with \eqref{predoubling} gives $0 < \lexp(\NN_\bp) \leq \uexp(\NN_\bp) < 2\delta$. Proposition \ref{propositiondoubling2} completes the proof.
\end{proof}

\begin{proof}[Proof of \text{(ii)}]
Since (up to equivalence) the only parabolic point of $G_n$ is the global fixed point of $H$ (Proposition \ref{propositionSproductGF}), the implication (A) \implies (B) follows from part (i). Conversely, suppose that (B) holds. Then by Proposition \ref{propositiondoubling2}, we have
\[
0 < \lexp(\NN_{\EE_\bp,H}) \leq \uexp(\NN_{\EE_\bp,H}) \leq 2\delta_n < \infty.
\]
Since $\delta_n\to \delta_H$ as $n\to\infty$ (Proposition \ref{propositionSproductexponent}(iv)), taking the limit and combining with the inequality $2\delta_H \leq \uexp(\NN_{\EE_\bp,H})$ yields \eqref{predoubling}. On the other hand, by Proposition \ref{propositiondoubling}, for each $n$, $G_n$ is of divergence type, so applying Proposition \ref{propositionSproductexponent}(iv) again, we see that $H$ is of divergence type.
%Assume that (A) is false, and we will show that (B) is false. We divide into cases according to which of the two conditions of Definition \ref{definitionpredoubling} fails:
%\begin{itemize}
%\item[Case 1.] Suppose that the inequalities $0 < \lexp(\NN_{\EE_\bp,H}) \leq \uexp(\NN_{\EE_\bp,H}) = 2\delta_H < \infty$ do not all hold. Since $\lexp(\NN_{\EE_\bp,H}) \leq \uexp(\NN_{\EE_\bp,H}) \leq 2\delta_H$, we must have either:
%\begin{itemize}
%\item[Case 1a.] $\lexp(\NN_{\EE_\bp,H}) = 0$. In this case, 
%\end{itemize}
%
%If $\lexp(\NN_{\EE_\bp,H}) = 0$, then by (I) of Proposition \ref{propositiondoubling}, none of the measures $\mu_n$ are doubling. Now the inequalities $\lexp(\NN_{\EE_\bp,H}) \leq 2\delta_H \leq \uexp(\NN_{\EE_\bp,H})$ and $\delta_H < \infty$ are automatic, so if (A) fails then we must have $\uexp(\NN_{\EE_\bp,H}) > 2\delta_H$. But then by choosing $n$ large enough, we may guarantee that $\delta_n < \uexp(\NN_{\EE_\bp,H})/2$ [more details needed\internal]. Then by Proposition \ref{propositiondoubling2}, $\mu_n$ is not doubling.
%\item[(II)] If $\uexp(\NN_{\EE_\bp,H}) > 2\delta_H$, then 
%\item[(III)] Suppose that $H$ is of convergence type. By choosing $n$ large enough, we may guarantee that $G_n$ is of convergence type [more details needed\internal]; then by (III) of Proposition \ref{propositiondoubling}, $\mu_n$ is not doubling.
%\end{itemize}
\end{proof}

\begin{example}
\label{exampledoubling}
If
\begin{equation}
\label{Nbppowerlaw}
\NN_\bp(R) \asymp_\times R^{2\delta_\bp} \all \bp\in P,
\end{equation}
then the groups $(G_\bp)_{\bp\in P}$ are pre-doubling, and thus by Proposition \ref{propositionpredoubling}(i), $\mu$ is doubling. Combining with Corollary \ref{corollaryGMFdoubling} gives
\[
\mu(B(\eta,e^{-t})) \asymp_\times e^{-\delta t} \begin{cases}
1 & \eta_t\notin \bigcup(\scrH)\\
e^{(2\delta_\xi - \delta) b(\eta,t)} & \eta_t\in H_\xi
\end{cases}.
\]
This generalizes B. Schapira's global measure formula \cite[Th\'eor\`eme 3.2]{Schapira} to the setting of regularly geodesic strongly hyperbolic metric spaces.
\end{example}

We remark that the asymptotic \eqref{Nbppowerlaw} is satisfied whenever $X$ is a finite-dimensional algebraic hyperbolic space; see e.g. \cite[Lemma 3.5]{Newberger}. In particular, specializing Schapira's global measure formula to the settings of finite-dimensional algebraic hyperbolic spacess and finite-dimensional real hyperbolic spaces give the global measure formulas of Newberger \cite[Main Theorem]{Newberger} and Stratmann--Velani--Sullivan \cite[Theorem 2]{StratmannVelani}, \cite[Theorem on p.271]{Sullivan_entropy}, respectively.

By contrast, when $X = \H = \H^\infty$, the asymptotic \eqref{Nbppowerlaw} is usually not satisfied. Let us summarize the various behaviors that we have seen for the orbital counting functions of parabolic groups acting on $\H^\infty$, and their implications for doubling:

\begin{examples}[Examples of doubling and non-doubling Patterson--Sullivan measures of geometrically finite subgroups of $\Isom(\H^\infty)$]
\label{exampledoublingnondoubling}
~
\begin{itemize}
\item[1.] In the proof of Theorem \ref{theoremnilpotentembedding} (cf. Remark \ref{remarknilpotentembedding}), we saw that if $\Gamma$ is a finitely generated virtually nilpotent group and if $f:\CO 1\infty\to\CO 1\infty$ is a function satisfying
\[
\alpha_\Gamma < \lexp(f) \leq \uexp(f) < \infty,
\]
then there exists a parabolic group $H\leq\Isom(\H^\infty)$ isomorphic to $\Gamma$ whose orbital counting function is asymptotic to $f$. Now, a group $H$ constructed in this way may or may not be pre-doubling; it depends on the chosen function $f$. We note that by applying Proposition \ref{propositionpredoubling}(ii) to such a group, one can construct examples of geometrically finite subgroups of $\Isom(\H^\infty)$ whose Patterson--Sullivan measures are not doubling. On the other hand, for any parabolic group $H$ constructed in this way, if $H$ is embedded into a geometrically finite group $G$ with sufficiently large Poincar\'e exponent (namely $2\delta_G > \uexp(f)$), then the Patterson--Sullivan measure of $G$ may be doubling (assuming that no other parabolic subgroups of $G$ are causing problems).
\item[2.] In Theorem \ref{theoremRtreeorbitalcounting}, we showed that if $f:\Rplus\to\N$ satisfies the condition
\[
\forall 0\leq R_1 \leq R_2 \; \text{$f(R_1)$ divides $f(R_2)$},
\]
then there exists a parabolic subgroup of $\Isom(\H^\infty)$ whose orbital counting function is equal to $f$. This provides even more examples of parabolic groups which are not pre-doubling. In particular, it provides examples of parabolic groups $H$ which satisfy either $\lexp(\NN_H) = 0$ or $\uexp(\NN_H) = \infty$ (cf. Example \ref{exampleparabolictorsion}); such groups cannot be embedded into any geometrically finite group with a doubling Patterson--Sullivan measure.
\end{itemize}
\end{examples}

Note that example 2 can be used to construct a geometrically finite group acting isometrically on an $\R$-tree which does not have a doubling Patterson--Sullivan measure. On the other hand, example 1 has no analogue in $\R$-trees by Remark \ref{remarkparabolicRtree}.

%Thus, there are a wealth of examples of both doubling and non-doubling Patterson--Sullivan measures of geometrically finite subgroups of $\Isom(\H^\infty)$.

%These examples indicate that the boundary between doubling and non-doubling behavior is quite subtle.% We end this section by noting that the constructions of parabolic groups mentioned above can be used to create examples of doubling measures as well as examples of non-doubling ones.



%\begin{observation}
%\label{observationhardyfield}
%Suppose that
%\[
%\NN_{\EE_\bp,H}(R) \asymp_\times \exp f \log(R)
%\]
%for some differentiable function $f:\Rplus\to\Rplus$ such that
%\[
%\lim_{t\to\infty} f'(t) \text{ exists.}
%\]
%Then $\dexp(\NN_{\EE_\bp,H}) = 2\delta_H = \lim_{t\to\infty} f'(t)$. In particular, $H$ is pre-doubling if and only if $H$ is of divergence type and $\delta_H > 0$.
%\end{observation}
%\begin{proof}
%We have
%\begin{align*}
%\dexp(\NN_{\EE_\bp,H}) = \dexp(\exp f\log)
%&= \lim_{\lambda,R\to\infty} \frac{f(\log(R) + \log(\lambda)) - f(\log(R))}{\log(\lambda)}\noreason\\
%&= \lim_{x,h\to\infty} \frac{f(x + h) - f(x)}{h} \noreason\\
%&= \lim_{t\to\infty} f'(t). \by{the fundamental theorem of calculus}
%\end{align*}
%Since $\dexp(\NN_{\EE_\bp,H})$ exists, $\dexp(\NN_{\EE_\bp,H}) = 2\delta_H$ by \eqref{poincarealternate}.
%\end{proof}


\section{Exact dimensionality of $\mu$}
\label{subsectionexactdimensional}
We now turn to the question of the fractal dimensions of the measure $\mu$. We recall that the \emph{Hausdorff dimension} and \emph{packing dimension} of a measure $\mu$ on $\del X$ are defined by the formulas
\begin{align*}
\HD(\mu) &= \inf\left\{\HD(A): \mu(\del X\butnot A) = 0\right\}\\
\PD(\mu) &= \inf\left\{\PD(A): \mu(\del X\butnot A) = 0\right\}.
\end{align*}
If $G$ is of convergence type, then $\mu$ is atomic, so $\HD(\mu) = \PD(\mu) = 0$. Consequently, for the remainder of this chapter we make the
\begin{standingassumption}
\label{assumptiondivergencetype}
$G$ is of divergence type.
\end{standingassumption}
Given this assumption, it is natural to expect that $\HD(\mu) = \PD(\mu) = \delta$. Indeed, the inequality $\HD(\mu)\leq\delta$ follows immediately from Theorems \ref{theorembishopjonesregular} and \ref{theoremGFcompact}, and in the Standard Case equality holds \cite[Proposiiton 4.10]{StratmannVelani}. Even stronger than the equalities $\HD(\mu) = \PD(\mu) = \delta$, it is natural to expect that $\mu$ is \emph{exact dimensional}:

\begin{definition}
\label{definitionexactdimensional}
A measure $\mu$ on a metric space $(Z,\Dist)$ is called \emph{exact dimensional of dimension $s$} if the limit
\begin{equation}
\label{dmudef}
d_\mu(\eta) := \lim_{t\to\infty} \frac1t \log\frac{1}{\mu(B(\eta,e^{-t}))}
\end{equation}
exists and equals $s$ for $\mu$-a.e. $\eta\in Z$.
\end{definition}

For example, every Ahlfors $s$-regular measure is exact dimensional of dimension $s$.

If the limit in \eqref{dmudef} does not exist, then we denote the lim inf by $\underline d_\mu(\eta)$ and the lim sup by $\overline d_\mu(\eta)$.

\begin{proposition}[{\cite[\68]{MSU}}]
\label{propositionexactdimensional}
For any measure $\mu$ on a metric space $(Z,\Dist)$,
\begin{align*}
\HD(\mu) &= \esssup_{\eta\in Z} \underline d_\mu(\eta)\\
\PD(\mu) &= \esssup_{\eta\in Z} \overline d_\mu(\eta).
\end{align*}
In particular, if $\mu$ is exact dimensional of dimension $s$, then
\[
\HD(\mu) = \PD(\mu) = s.
\]
\end{proposition}

Combining Proposition \ref{propositionexactdimensional} with Lemma \ref{lemmaCCBglobalmeasure} and Observation \ref{observationnonatomic} immediately yields the following:

\begin{observation}
\label{observationHDdeltaPD}
If $\mu$ is the Patterson--Sullivan measure of a geometrically finite group of divergence type, then
\[
\HD(\mu) \leq \delta \leq \PD(\mu).
\]
In particular, if $\mu$ is exact dimensional, then $\mu$ is exact dimensional of dimension $\delta$.
\end{observation}

It turns out that $\mu$ is not necessarily exact dimensional (Example \ref{examplenotexactdim}), but counterexamples to exact dimensionality must fall within a very narrow window (Theorem \ref{theoremhlogseries}), and in particular if $\mu$ is doubling then $\mu$ is exact dimensional (Corollary \ref{corollarydoublingexactdim}). As a first step towards these results, we will show that exact dimensionality is equivalent to a certain Diophantine condition. For this, we need to recall some results from \cite{FSU4}.

\subsection{Diophantine approximation on $\Lambda$}
\label{subsubsectiondiophantineapproximation}
Classically, Diophantine approximation is concerned with the approximation of a point $x\in\R\butnot\Q$ by a rational number $p/q\in\Q$. The two important quantities are the \emph{error term} $|x - p/q|$ and the \emph{height} $q$. Given a function $\Psi:\N\to\Rplus$, the point $x\in\R\butnot\Q$ is said to be \emph{$\Psi$-approximable} if
\[
\left|x - \frac pq\right| \leq \Psi(q) \text{ for infinitely many $p/q\in\Q$}.
\]
In the setting of a group acting on a hyperbolic metric space, we can instead talk about \emph{dynamical} Diophantine approximation, which is concerned with the approximation of a point $\eta\in\Lambda$ by points $g(\xi)\in G(\xi)$, where $\xi\in\Lambda$ is a distinguished point. For this to make sense, one needs a new definition of error and height: the error term is defined to be $\Dist(g(\xi),\eta)$, and the height is defined to be $b^{\dogo g}$. (If there is more than one possibility for $g$, it may be chosen so as to minimize the height.) Some motivation for these definitions comes from considering classical Diophantine approximation as a special case of dynamical Diophantine approximation which occurs when $X = \H^2$ and $G = \SL_2(\Z)$; see e.g. \cite[Observation 1.15]{FSU4} for more details. Given a function $\Phi:\Rplus\to (0,\infty)$, the point $\eta\in\Lambda$ is said to be \emph{$\Phi,\xi$-well approximable} if for every $K > 0$ there exists $g\in G$ such that
\[
\Dist(g(\xi),\eta) \leq \Phi(K b^{\dogo g}) \text{ for infinitely many $g\in G$}
\]
(cf. \cite[Definition 1.36]{FSU4}). Moreover, $\eta$ is said to be \emph{$\xi$-very well approximable} if
\[
\omega_\xi(\eta) := \limsup_{\substack{g\in G \\ g(\xi)\to\eta}} \frac{-\log_b \Dist(g(\xi),\eta)}{\dogo g} > 1
\]
%for every $\epsilon > 0$ we have
%\[
%\Dist(g(\xi),\eta) \leq b^{-(1 + \epsilon)\dogo g} \text{ for infinitely many $g\in G$}
%\]
(cf. \cite[p.9]{FSU4}). The set of $\Phi,\xi$-well approximable points is denoted $\WA_{\Phi,\xi}$, while the set of $\xi$-very well approximable points is denoted $\VWA_\xi$. Finally, a point $\eta$ is said to be \emph{Liouville} if $\omega_\xi(\eta) = \infty$; the set of Liouville points is denoted $\Liou_\xi$.

In the following theorems, we return to the setting of Standing Assumptions \ref{standingassumptionsGFmeasures} and \ref{assumptiondivergencetype}.

\begin{theorem}[Corollary of {\cite[Theorem 8.1]{FSU4}}]
\label{theoremFSUkhinchin}
Fix $\bp\in P$, and let $\Phi:\Rplus\to(0,\infty)$ be a function such that the function $t\mapsto t\Phi(t)$ is nonincreasing. Then
\begin{itemize}
\item[(i)] $\mu(\WA_{\Phi,\bp}) = 0\text{ or }1$ according to whether the series
\begin{equation}
\label{WAseries}
\sum_{g\in G} e^{-\delta \dogo g} \II_\bp\left(\frac{1}{e^{\dogo g} \Phi(K e^{\dogo g})}\right)
\end{equation}
converges for some $K > 0$ or diverges for all $K > 0$, respectively.
\item[(ii)] $\mu(\VWA_\bp) = 0\text{ or }1$ according to whether the series
\begin{equation}
\label{VWAseries}
\SigmaVWA(\bp,\kappa) := \sum_{g\in G} e^{-\delta \dogo g} \II_\bp(e^{\kappa\dogo g})
\end{equation}
converges for all $\kappa > 0$ or diverges for some $\kappa > 0$, respectively.
\item[(iii)] $\mu(\Liou_\bp) = 0\text{ or }1$ according to whether the series $\SigmaVWA(\bp,\kappa)$ converges for some $\kappa > 0$ or diverges for all $\kappa > 0$, respectively.
\end{itemize}
\end{theorem}
\begin{proof}
Standing Assumption \ref{assumptiondivergencetype}, Theorem \ref{theoremahlforsthurstongeneral}, and Observation \ref{observationnonatomic} imply that $\mu$ is ergodic and that $\mu(\bp) = 0$, thus verifying the hypotheses of \cite[Theorem 8.1]{FSU4}. Theorem \ref{theoremglobalmeasure} shows that
\[
\II_\bp(C_1/r) \lesssim_{\times,\bp} \mu(B(\bp,r)) \lesssim_{\times,\bp} \II_\bp(C_2/r)
\]
for some constants $C_1 \geq 1 \geq C_2 > 0$. Thus for all $K > 0$,
\begin{align*}
&\sum_{g\in G} e^{-\delta\dogo g} \II_\bp\left(\frac{1}{e^{\dogo g} \Phi(K C_1 e^{\dogo g})}\right)\\
&\leq_\pt \sum_{g\in G} e^{-\delta\dogo g} \II_\bp\left(\frac{C_1}{e^{\dogo g} \Phi(K e^{\dogo g})}\right)\\
&\lesssim_\times \text{\cite[(8.1)]{FSU4}}\\
&\lesssim_\times \sum_{g\in G} e^{-\delta\dogo g} \II_\bp\left(\frac{C_2}{e^{\dogo g} \Phi(K e^{\dogo g})}\right)\\
&\leq_\pt \sum_{g\in G} e^{-\delta\dogo g} \II_\bp\left(\frac{1}{e^{\dogo g} \Phi((K/C_1) e^{\dogo g})}\right).
\end{align*}
Thus, \cite[(8.1)]{FSU4} diverges for all $K > 0$ if and only if \eqref{WAseries} diverges for all $K > 0$. This completes the proof of (i). To demonstrate (ii) and (iii), simply note that $\VWA_\bp = \bigcup_{c > 0} \WA_{\Phi_c,\bp}$ and $\Liou_\bp = \bigcap_{c > 0} \WA_{\Phi_c,\bp}$, where $\Phi_c(t) = t^{-(1 + c)}$, and apply (i). The constant $K$ may be absorbed by a slight change of $\kappa$.
\end{proof}

\begin{theorem}[Corollary of {\cite[Theorem 7.1]{FSU4}}]
\label{theoremFSUjarnik}
For all $\xi\in\Lambda$ and $c > 0$,
\[
\HD(\WA_{\Phi_c,\xi}) \leq \frac{\delta}{1 + c},
\]
where $\Phi_c(t) = t^{-(1 + c)}$ as above. In particular, $\HD(\Liou_\xi) = 0$, and $\VWA_\xi$ can be written as the countable union of sets of Hausdorff dimension strictly less than $\delta$.
\end{theorem}
\noindent (No proof is needed as this follows directly from \cite[Theorem 7.1]{FSU4}.)

There is a relation between dynamical Diophantine approximation by the orbits of parabolic points and the lengths of cusp excursions along geodesics. A well-known example is that a point $\eta\in\Lambda$ is dynamically badly approximable with respect to every parabolic point if and only if the geodesic $\geo\zero\eta$ has bounded cusp excursion lengths \cite[Proposition 1.21]{FSU4}. The following observation is in a similar vein:

\begin{observation}
\label{observationVWAbetat}
For $\eta\in\Lambda$, we have:
\begin{align*}
\eta\in\bigcup_{p\in P}\VWA_\bp
&\;\;\Leftrightarrow\;\; \limsup_{\substack{\xi\in\Lbp \\ t_\xi\to\infty}} \frac{\lb\xi|\eta\rb - t_\xi}{t_\xi} > 0
\;\;\Leftrightarrow\;\; \limsup_{t\to\infty} \frac{b(\eta,t)}{t} > 0\\
\eta\in\bigcup_{p\in P}\Liou_\bp
&\;\;\Leftrightarrow\;\; \limsup_{\substack{\xi\in\Lbp \\ t_\xi\to\infty}} \frac{\lb\xi|\eta\rb - t_\xi}{t_\xi} = \infty
\;\;\Leftrightarrow\;\; \limsup_{t\to\infty} \frac{b(\eta,t)}{t} = 1.
\end{align*}
%\begin{align} \tag{A}
%&\eta\notin\bigcup_{p\in P}\VWA_\bp\\ \tag{B}
%&\limsup_{\substack{\xi\in\Lbp \\ t_\xi\to\infty}} \frac{\lb\xi|\eta\rb - t_\xi}{t_\xi} = 0 \\ \tag{C}
%&\limsup_{t\to\infty} \frac{b(\eta,t)}{t} = 0.
%\end{align}
\end{observation}
\begin{proof}
If $\xi = g(\bp)\in\Lbp$, then $\dogo g \gtrsim_\plus t_\xi$, with $\asymp_\plus$ for at least one value of $g$ (Lemma \ref{lemmatop}). Thus
\begin{align*}
\max_{\bp\in P}\omega_\bp(\eta) = \max_{\bp\in P}\limsup_{\substack{g\in G \\ g(\bp)\to\eta}} \frac{\log\Dist(g(\bp),\eta)}{\dogo g}
= \limsup_{\substack{\xi\in\Lbp \\ \xi\to\eta}} \frac{\lb \xi|\eta\rb}{t_\xi},
\end{align*}
so
\begin{equation}
\label{VWAbetat1}
\limsup_{\substack{\xi\in\Lbp \\ t_\xi\to\infty}} \frac{\lb\xi|\eta\rb - t_\xi}{t_\xi} = \max_{\bp\in P}\omega_\bp(\eta) - 1.
%\begin{cases}
%= \infty & \eta\in\bigcup_{p\in P}\Liou_\bp\\
%\in (0,\infty) & \eta\in\bigcup_{p\in P}\VWA_\bp\butnot \bigcup_{p\in P}\Liou_\bp\\
%= 0 & \eta\notin\bigcup_{p\in P}\VWA_\bp
%\end{cases}.
\end{equation}
On the other hand, it is readily verified that if $\geo\zero\eta$ intersects $H_\xi$, then the function $f(t) = b(\eta,t)/t$ attains its maximum at $t = \lb \xi|\eta\rb_\zero$, at which $f(t) = \lb \xi|\eta\rb_\zero - t_\xi$. Thus we have that
\begin{equation}
\label{VWAbetat2}
\begin{split}
\limsup_{t\to\infty} \frac{b(\eta,t)}{t}
= \limsup_{\substack{\xi\in\Lbp \\ t_\xi\to\infty}} \sup_{\substack{t > 0 \\ \eta_t\in H_\xi}} \frac{b(\eta,t)}{t}
&= \limsup_{\substack{\xi\in\Lbp \\ t_\xi\to\infty}} \frac{\lb \xi|\eta\rb_\zero - t_\xi}{\lb \xi|\eta\rb_\zero}\\
&= 1 - \frac{1}{\max_{\bp\in P}\omega_\bp(\eta)}
\end{split}
%\begin{cases}
%= 1 & \eta\in\bigcup_{p\in P}\Liou_\bp\\
%\in (0,1) & \eta\in\bigcup_{p\in P}\VWA_\bp\butnot \bigcup_{p\in P}\Liou_\bp\\
%= 0 & \eta\notin\bigcup_{p\in P}\VWA_\bp
%\end{cases}.
\end{equation}
Since
\[
\max_{\bp\in P}\omega_\bp(\eta) \begin{cases}
= \infty & \eta\in\bigcup_{p\in P}\Liou_\bp\\
\in (1,\infty) & \eta\in\bigcup_{p\in P}\VWA_\bp\butnot \bigcup_{p\in P}\Liou_\bp\\
= 1 & \eta\notin\bigcup_{p\in P}\VWA_\bp
\end{cases}
\]
applying \eqref{VWAbetat1} and \eqref{VWAbetat2} completes the proof.
\end{proof}

We are now ready to state our main theorem regarding the relation between exact dimensionality and dynamical Diophantine approximation:

\begin{theorem}
\label{theoremequivalentVWA}
The following are equivalent:
\begin{itemize}
\item[(A)] $\mu(\VWA_\bp) = 0\all \bp\in P$.
\item[(B)] $\mu$ is exact dimensional.
\item[(C)] $\HD(\mu) = \delta$.
\item[(D)] $\mu(\VWA_\xi) = 0\all \xi\in\Lambda$.
\end{itemize}
\end{theorem}
The implication (B) \implies (C) is part of Proposition \ref{propositionexactdimensional}, while (C) \implies (D) is an immediate consequence of Theorem \ref{theoremFSUjarnik}, and (D) \implies (A) is trivial. Thus we demonstrate (A) \implies (B):

\begin{proof}[Proof of \text{(A) \implies (B)}]
Fix $\eta\in\Lambda\butnot\bigcup_{\bp\in P}\VWA_\bp$ and $t > 0$. Suppose that $\eta_t\in H_\xi$ for some $\xi\in\Lbp$. Let $t_- < t < t_+$ satisfy
\[
t_- \asymp_\plus t_\xi, \; t_+ \asymp_\plus 2\lb \xi|\eta\rb_\zero - t_\xi, \text{ and } \eta_{t_\pm} \notin\bigcup(\scrH).
\]
Then by Lemma \ref{lemmaCCBglobalmeasure},
\[
\mu(B(\eta,e^{-t_\pm})) \asymp_\times e^{-\delta t_\pm}.
\]
In particular
\begin{equation}
\label{dt-dt+bounds}
\delta t_- \lesssim_\plus \log\frac{1}{\mu(B(\eta,e^{-t}))} \lesssim_\plus \delta t_+.
\end{equation}
Now, by Observation \ref{observationVWAbetat}, we have
\[
\frac{t_+ - t_-}{t} \leq \frac{2(\lb \xi|\eta\rb_\zero - t_\xi + (\text{constant}))}{t_\xi} \to 0 \text{ as } t\to\infty.
\]
Since $t_- < t < t_+$, it follows that $t_-/t,t_+/t\to 1$ as $t\to \infty$. Combining with \eqref{dt-dt+bounds} gives $d_\mu(\eta) = \delta$ (cf. \eqref{dmudef}). But by assumption (A), this is true for $\mu$-a.e. $\eta\in\Lambda$. Thus $\mu$ is exact dimensional.
%It follows that
%\[
%d_\mu(\eta) := \lim_{t\to\infty} \frac1t \log\frac{1}{\mu(B(\eta,e^{-t}))} = \delta,
%\]
%i.e. the pointwise dimension of $\mu$ at $\eta$ is equal to $\delta$.
\end{proof}

%old version:
%\subsection{Examples of exact dimensional and non-exact dimensional measures}
%new shortened version:
\subsection{Examples and non-examples of exact dimensional measures}
Combining Theorems \ref{theoremequivalentVWA} and \ref{theoremFSUkhinchin} gives a necessary and sufficient condition for $\mu$ to be exact dimensional in terms of the convergence or divergence of a family of series. We can ask how often this condition is satisfied. Our first result shows that it is almost always satisfied:

\begin{theorem}
\label{theoremhlogseries}
If for all $\bp\in P$, the series
\begin{equation}
\label{hlogseries}
\sum_{h\in G_\bp} e^{-\delta\dogo h} \dogo h
\asymp_\times \sum_{h\in G_\bp} \|h\|_\bp^{-2\delta} \log\|h\|_\bp
\asymp_\times \sum_{k = 0}^\infty e^{-2\delta k} k \NN_\bp(e^k)
\end{equation}
converges, then $\mu$ is exact dimensional.
\end{theorem}
\begin{proof}
Fix $\bp\in P$ and $\kappa > 0$. We have
\begin{align*}
\SigmaVWA(\bp,\kappa)
&=_\pt \sum_{g\in G} e^{-\delta\dogo g} \sum_{\substack{h\in G_\bp \\ \dogo h \geq \kappa\dogo g/2}} e^{-\delta\dogo h}\\
&=_\pt \sum_{h\in G_\bp} e^{-\delta\dogo h} \sum_{\substack{g\in G \\ \dogo g \leq 2\dogo h/\kappa}} e^{-\delta\dogo g}\\
&\asymp_\times \sum_{h\in G_\bp} e^{-\delta\dogo h} \sum_{k\leq 2\dogo h/\kappa + 1} e^{-\delta k} \#\{g\in G : k - 1 \leq \dogo g < k\} \noreason\\
&\leq_\pt \sum_{h\in G_\bp} e^{-\delta\dogo h} \sum_{k\leq 2\dogo h/\kappa + 1} e^{-\delta k} \NN_{X,G}(k) \noreason\\
&\lesssim_\times \sum_{h\in G_\bp} e^{-\delta\dogo h} \sum_{k\leq 2\dogo h/\kappa + 1} 1 \by{Corollary \ref{corollaryupperorbitalbound}}\\
&\asymp_\times \sum_{h\in G_\bp} e^{-\delta\dogo h} \dogo h.
\end{align*}
% &\lesssim_\times \sum_{g\in G} e^{-\delta\dogo g} e^{-2\delta\kappa\dogo g} \w\II(e^{\kappa\dogo g})\\
%&=_\pt \sum_{g\in G} \sum_{k = 1}^\infty e^{-\delta\dogo g} e^{-2\delta(\kappa\dogo g + k)} \NN_\bp(e^{\kappa\dogo g + k})\\
So if \eqref{hlogseries} converges, so does $\SigmaVWA(\bp,\kappa)$, and thus by Theorems \ref{theoremFSUkhinchin} and \ref{theoremequivalentVWA}, $\mu$ is exact dimensional.
\end{proof}

\begin{corollary}
\label{corollarydeltagtrdeltap}
If for all $\bp\in P$, $\delta_\bp < \delta$, then $\mu$ is exact dimensional.
\end{corollary}
\begin{proof}
In this case, the series \eqref{hlogseries} converges, as it is dominated by $\Sigma_s(G_\bp)$ for any $s\in (\delta_\bp,\delta)$.
\end{proof}

\begin{remark}
\label{remarkdeltagtrdeltap}
Combining with Proposition \ref{propositiondeltagtrdeltap} shows that if $\mu$ is not exact dimensional, then
\[
\sum_{h\in G_\bp} e^{-\delta\dogo h} < \infty = \sum_{h\in G_\bp} e^{-\delta\dogo h} \dogo h
\]
for some $\bp\in P$. Equivalently,
\[
\sum_{k = 0}^\infty e^{-2\delta k} \NN_\bp(e^k) < \infty = \sum_{k = 0}^\infty e^{-2\delta k} k \NN_\bp(e^k).
\]
This creates a very ``narrow window'' for the orbital counting function $\NN_\bp$.
\end{remark}

\begin{corollary}
\label{corollarydoublingexactdim}
If $\mu$ is doubling, then $\mu$ is exact dimensional.
\end{corollary}
\begin{proof}
If $\mu$ is doubling, then
\begin{align*}
\sum_{k = 0}^\infty e^{-2\delta k} k \NN_\bp(e^k)
&=_\pt \sum_{k = 1}^\infty \sum_{\ell = 0}^\infty e^{-2\delta (k + \ell)} \NN_\bp(e^{k + \ell})\\
&=_\pt \sum_{k = 1}^\infty e^{-2\delta k} \w\II_\bp(e^k)\\
&\asymp_\times \sum_{k = 1}^\infty e^{-2\delta k} \NN_\bp(e^k). \by{Proposition \ref{propositiondoubling}}
\end{align*}
Remark \ref{remarkdeltagtrdeltap} completes the proof.
\end{proof}


Our next theorem shows that in certain circumstances, the converse holds in Theorem \ref{theoremhlogseries}. Specifically:

\begin{theorem}
\label{theoremhlogseriesconverse}
Suppose that $X$ is an $\R$-tree and that $G$ is the pure Schottky product (cf. Definition \ref{definitionschottkyproductRtree}) of a parabolic group $H$ with a lineal group $J$. Let $\bp$ be the global fixed point of $H$, so that $P = \{\bp\}$ is a complete set of inequivalent parabolic points for $G$ (Proposition \ref{propositionSproductGF}). Suppose that the series \eqref{hlogseries} diverges. Then $\mu$ is not exact dimensional; moreover, $\mu(\Liou_\bp) = 1$ and $\dim_H(\mu) = 0$.
\end{theorem}
\begin{example}
\label{examplenotexactdim}
To see that the hypotheses of this theorem are not vacuous, fix $\delta > 0$ and let
\[
f(R) = \frac{R^{2\delta}}{\log^2(R)},
\]
or more generally, let $f$ be any increasing function such that $\sum_1^\infty e^{-2\delta k} k f(e^k)$ diverges but $\sum_1^\infty e^{-2\delta k} f(e^k)$ converges. By Theorem \ref{theoremRtreeorbitalcounting}, there exists an $\R$-tree $X$ and a parabolic group $H\leq\Isom(X)$ such that $\NN_{\EE_\bp,H} \asymp_\times f$. Then the series \eqref{hlogseries} diverges, but $\Sigma_\delta(H) < \infty$. Thus, there exists a unique $r > 0$ such that
\[
2\sum_{n = 1}^\infty e^{-\delta r} = \frac{1}{\Sigma_\delta(H) - 1}\cdot
\]
Let $J = r\Z$, interpreted as a group acting by translations on the $\R$-tree $\R$, and let $G$ be the pure Schottky product of $H$ and $J$. Then $\Sigma_\delta(J) - 1 = 2\sum_{n = 1}^\infty e^{-\delta r}$, so $(\Sigma_\delta(H) - 1)(\Sigma_\delta(J) - 1) = 1$. Since the map $s\mapsto (\Sigma_s(H) - 1)(\Sigma_s(J) - 1)$ is decreasing, it follows from Proposition \ref{propositionschottkyproductRtree} that $\Delta(G) = [0,\delta]$. In particular, $G$ is of divergence type, so Standing Assumption \ref{assumptiondivergencetype} is satisfied.
\end{example}
\begin{remark}
Applying a BIM embedding allows us to construct an example acting on $\H^\infty$.
\end{remark}
\begin{proof}[Proof of Theorem \ref{theoremhlogseriesconverse}]
As in the proof of Proposition \ref{propositionschottkyproductRtree} we let
\[
E = (H\butnot\{\id\})(J\butnot\{\id\}),
\]
so that
\[
G = \bigcup_{n\geq 0} J E^n H.% = \bigcup_{n = 0}^\infty \bigcup_{j_0\in J} \bigcup_{\substack{h_1,\ldots,h_n\in H^* \\ j_1,\ldots,j_n \in J^*}} \bigcup_{h_{n + 1}\in H} \{j_0 h_1 j_1 \cdots h_n j_n h_{n + 1}\}.
\]
Define a measure $\theta$ on $E$ via the formula
%Then by \eqref{schottkyproductRtree}, we have for all $s\geq 0$
%\begin{align*}
%\Sigma_s(G)
%= \sum_{g\in G} e^{-s\dogo g}
%&= \sum_{n = 0}^\infty \sum_{j_0\in J} \sum_{g_1,\ldots,g_n\in E} \sum_{h_{n + 1}\in H} e^{-s[\|j_0\| + \sum_1^n \|g_i\| + \|h_{n + 1}\|]}\\
%&= \Sigma_s(J) \Sigma_s(H) \sum_{n = 0}^\infty \left(\sum_{g\in E} e^{-s\|g\|}\right)^n.
%\end{align*}
%Combining with the definition of $\delta$, the fact that $G$ is of divergence type (Standing Assumption \ref{assumptiondivergencetype}), and the fact that $\Sigma_\delta(J),\Sigma_\delta(H) < \infty$ (Proposition \ref{propositiondeltagtrdeltap}), we see that $\sum_{g\in E} e^{-\delta\|g\|} = 1$, so the measure $\mu$ on $E$ defined by the formula
\[
\theta = \sum_{g\in E} e^{-\delta\dogo g} \delta_g.
\]
By Proposition \ref{propositionschottkyproductRtree}, the fact that $G$ is of divergence type (Standing Assumption \ref{assumptiondivergencetype}), and the fact that $\Sigma_\delta(J),\Sigma_\delta(H) < \infty$ (Proposition \ref{propositiondeltagtrdeltap}), $\theta$ is a probability measure. The Patterson--Sullivan measure of $G$ is related to $\theta$ by the formula
\[
\mu = \frac{1}{\Sigma_\delta(J) - 1}\sum_{j\in J} e^{-\delta\dogo j} j_* \pi_*[\theta^\N],
\]
where $\pi:E^\N \to \Lambda_G$ is the coding map.

Next, we use a theorem proven independently by H. Kesten and A. Raugi,\Footnote{We are grateful to ``cardinal'' of \url{http://mathoverflow.net} and J. P. Conze, respectively, for these references.} which we rephrase here in the language of measure theory:

\begin{theorem}[\cite{Kesten}; see also \cite{Raugi}]
\label{theoremkesten}
Let $\theta$ be a probability measure on a set $E$, and let $f:E\to\R$ be a function such that
\[
\int |f(x)| \;\dee\theta(x) = \infty.
\]
Then for $\theta^\N$-a.e. $(x_n)_1^\infty \in \R^\N$,
\[
\limsup_{n\to\infty} \frac{|f(x_{n + 1})|}{|\sum_1^n f(x_i)|} = \infty.
\]
\end{theorem}

Letting $f(g) = \dogo g$, the theorem applies to our measure $\theta$, because our assumption that \eqref{hlogseries} diverges is equivalent to the assertion that $\int f(x) \;\dee\theta(x) = \infty$. Now fix $j\in J$, and let $(g_n)_1^\infty \in E^\N$ be a $\theta^\N$-typical point. Then the limit point
\[
\eta = \lim_{n\to\infty} j g_1\cdots g_n(\zero)
\]
represents a typical point with respect to the Patterson--Sullivan measure $\mu$. By Theorem \ref{theoremkesten}, we have
\[
\limsup_{n\to\infty} \frac{\dogo{g_{n + 1}}}{\sum_1^n \dogo{g_i}} = \infty.
\]
Write $g_i = h_i j_i$ for each $i$. Then $\dogo{g_i} = \dogo{h_i} + \dogo{j_i}$. Since $\int \dogo{j}\dee\theta(hj) < \infty$, the law of large numbers implies that $\lim_{n \to \infty} \frac{\dogo{j_{n + 1}}}{\sum_1^n \dogo{j_i}} < \infty$, so
\[
\limsup_{n\to\infty} \frac{\dogo{h_{n + 1}}}{\sum_1^n \dogo{g_i}} = \infty.
\]
But $\dogo{h_{n + 1}}$ represents the length of the excursion of the geodesic $\geo\zero\eta$ into the cusp corresponding to the parabolic point $g_1\cdots g_n(\bp)$. Combining with Observation \ref{observationVWAbetat} shows that $\eta\in\Liou_\bp$. Since $\eta$ was a $\mu$-typical point, this shows that $\mu(\Liou_\bp) = 1$. By Theorem \ref{theoremFSUjarnik}, this implies that $\HD(\mu) = 0$. By Observation \ref{observationHDdeltaPD}, $\mu$ is not exact dimensional.
\end{proof}

\endinput
